\chapter{Theorem B by a sequential preparation protocol}
\label{ch:one_dimensional_protocol}

This chapter gives the proof of Theorem~B in
\cite[Part~II]{TNLean2026OneDimensional}, which specializes the preparation
argument of \cite{OpenAI2026PolynomialPEPS} for two-dimensional lattices,
formalized in
Chapters~\ref{ch:peps_approximation}--\ref{ch:peps_dyadic_protocol}, to the
chain. Local circuits sweep across the chain in stages and prepare the
ground state block by block. The information estimates of
Chapter~\ref{ch:one_dimensional_chain} bound the error made at each stage,
and the structure of the circuits bounds the Schmidt rank across every cut.
Theorem~\ref{thm:chain_protocol_mps_approximation} is the resulting
approximation, within distance $1/L$ of the ground state up to a global
phase. Chapter~\ref{ch:one_dimensional_chain} proves Theorem~B by sequential
truncation instead.

\section{Physical coordinates of the terminal chain registers}

\begin{definition}[Terminal main and auxiliary registers]
    \label{def:chain_terminal_registers}
    \lean{
        TNLean.MPS.AreaLaw.ChainSweep.stageCount,
        TNLean.MPS.AreaLaw.ChainSweep.Party,
        TNLean.MPS.AreaLaw.ChainSweep.AuxiliaryCopy,
        TNLean.MPS.AreaLaw.ChainSweep.AuxiliarySite,
        TNLean.MPS.AreaLaw.ChainSweep.RawLabel,
        TNLean.MPS.AreaLaw.ChainSweep.mainOwner,
        TNLean.MPS.AreaLaw.ChainSweep.auxiliaryOwner,
        TNLean.MPS.AreaLaw.ChainSweep.rawOwner,
        TNLean.MPS.AreaLaw.ChainSweep.MainSites,
        TNLean.MPS.AreaLaw.ChainSweep.AuxiliarySites,
        TNLean.MPS.AreaLaw.ChainSweep.physicalDimension
    }
    \leanok
    Fix integers $L,q,w\ge1$ with $L\ge60w$, put $t=20w$ and
    $m=\lfloor L/t\rfloor-3$, and take the distinct parties
    $R_0,\ldots,R_m,P_1,\ldots,P_m$.
    The main sites are $X=\{0,\ldots,L-1\}$. The owner of $x<mt$ is
    $P_{\lfloor x/t\rfloor+1}$, and the owner of $x\ge mt$ is $R_m$.
    There is a right auxiliary copy at each stage $0\le k<m$ and a left
    auxiliary copy at each stage $1\le k<m$; every copy has all sites of
    $X$. Equal centres do not identify these distinct labelled copies.
    For a right copy let $a=(k+1)t$ and $(A,B,C)=(R_k,R_{k+1},P_{k+1})$;
    for a left copy let $a=kt$ and $(A,B,C)=(R_k,P_{k+1},P_k)$.
    Its auxiliary site $x$ is owned by $A$ on $[a-3w,a-2w)$, by $B$ on
    $[a-2w,a-w)$ and by $C$ elsewhere. Write $\mathcal A$ for the set of
    all auxiliary copy-site labels, and $M_p\subseteq X$ and
    $\mathcal A_p\subseteq\mathcal A$ for the main and auxiliary labels
    at party $p$.
    The physical factor at $p$ has dimension $d_p=q^{|M_p|}$, whereas
    every auxiliary factor remains private at its specified owner.
    The distinction between these factors follows the corrected terminal
    assertion in \cite{gap:tnlean_one_dimensional_band_sweep_corrections}.
\end{definition}

\begin{lemma}[Terminal block sizes and party counts]
    \label{thm:chain_terminal_block_sizes}
    \lean{
        TNLean.MPS.AreaLaw.ChainSweep.terminal_length_eq,
        TNLean.MPS.AreaLaw.ChainSweep.terminal_length_bounds,
        TNLean.MPS.AreaLaw.ChainSweep.card_mainSites_inr_le,
        TNLean.MPS.AreaLaw.ChainSweep.card_mainSites_inl_lt,
        TNLean.MPS.AreaLaw.ChainSweep.twenty_mul_stageCount_le,
        TNLean.MPS.AreaLaw.ChainSweep.card_party_le,
        TNLean.MPS.AreaLaw.ChainSweep.physicalDimension_le
    }
    \leanok
    \uses{def:chain_terminal_registers}
    The final suffix has length in $[3t,4t)$, every completed main
    block has at most $t$ sites, and every remainder party has fewer
    than $4t$ main sites. In particular $d_p\le q^{4t}$ for every party.
    Moreover $20m\le L$ and the number of parties is $2m+1\le2L$.
\end{lemma}
\begin{proof}\leanok
    \uses{def:chain_terminal_registers}
    Euclidean division gives $L=(m+3)t+r$ with $0\le r<t$.
    On a fixed completed main fibre, $x\mapsto x\bmod t$ is injective,
    since $\lfloor x/t\rfloor$ is fixed. On every remainder main fibre,
    $x\mapsto x-mt$ is injective into the final suffix. These injections
    give the cardinal bounds, and monotonicity of powers for $q\ge1$
    gives the dimension estimate. Since $t\ge20$ and $mt\le L$, one
    has $20m\le L$; the party-count bound follows.
\end{proof}

\begin{lemma}[Margins at the sweep stages]
    \label{thm:chain_stage_margins}
    \lean{
        TNLean.MPS.AreaLaw.ChainSweep.stage_right_margin,
        TNLean.MPS.AreaLaw.ChainSweep.stage_left_margin
    }
    \leanok
    \uses{def:chain_terminal_registers}
    For every stage $0\le k<m$, one has $(k+4)t\le L$.
    If $k\ge1$, one also has $8w\le kt$.
\end{lemma}
\begin{proof}
    \leanok
    \uses{thm:chain_terminal_block_sizes}
    Since $k+1\le m$ and $L=mt+3t+r$ with $r\ge0$,
    $(k+4)t\le mt+3t\le L$. For $k\ge1$,
    $kt\ge t=20w\ge8w$.
\end{proof}

\begin{lemma}[Private preparation of the initial copies]
    \label{thm:chain_initial_copies}
    \lean{
        TNLean.MPS.AreaLaw.ChainSweep.initialAuxiliaryOrder,
        TNLean.MPS.AreaLaw.ChainSweep.initialAuxiliaryHolder,
        TNLean.MPS.AreaLaw.ChainSweep.initialFrames,
        TNLean.MPS.AreaLaw.ChainSweep.length_initialFrames,
        TNLean.MPS.AreaLaw.ChainSweep.mem_initialFrames,
        TNLean.MPS.AreaLaw.ChainSweep.norm_initialFrames,
        TNLean.MPS.AreaLaw.ChainSweep.exists_initialization
    }
    \leanok
    \uses{def:chain_terminal_registers}
    Let $\Omega$ be a unit vector on the chain. Initially place the main
    copy at $R_0$ and each auxiliary copy at its specified party $C$.
    There are $1+m+\max\{m-1,0\}$ copies, each held entirely by one
    party. Take a party's private memory to be the tensor product of
    the copies it holds, with their product reference vector. These
    vectors have norm one. Private maps and changes of tensor-factor
    order prepare the complete copy reference exactly, with no
    nonprivate operation.
\end{lemma}
\begin{proof}
    \leanok
    \uses{def:chain_terminal_registers, thm:lattice_sheet_initialization}
    Each right copy at stage $k$ is held by $P_{k+1}$, and each left
    copy at stage $k+1$ is held by $P_{k+1}$. The copies remain
    distinct even when their holders coincide. Grouping the copies
    by their holders gives a product across parties; each grouped
    vector has norm equal to the product of the copy norms, hence one.
    Theorem~\ref{thm:lattice_sheet_initialization} splits each party's
    private memory into its labelled copies and restores their prescribed
    order. Every map is private, and its action on the product vector
    is exact.
\end{proof}

\begin{definition}[Block coordinates]
    \label{def:chain_terminal_coordinates}
    \uses{def:chain_terminal_registers}
    \lean{
        TNLean.MPS.AreaLaw.ChainSweep.blockEncoding,
        TNLean.MPS.AreaLaw.ChainSweep.mainConfigurationEquiv,
        TNLean.MPS.AreaLaw.ChainSweep.terminalConfigurationEquiv,
        TNLean.MPS.AreaLaw.ChainSweep.localConfigurationEquiv,
        TNLean.MPS.AreaLaw.ChainSweep.privateLayout,
        TNLean.MPS.AreaLaw.ChainSweep.physicalOrder,
        TNLean.MPS.AreaLaw.ChainSweep.sheetReadout,
        TNLean.MPS.AreaLaw.ChainSweep.terminalRegrouping
    }
    \leanok
    At every party $p$, use increasing main-site order to identify
    configurations on $M_p$ with the basis of $\mathbb C^{d_p}$.
    Denote this identification by $\chi_p$, and put
    $\chi(c)_p=\chi_p(c|_{M_p})$ for a main configuration $c$.
    The terminal coordinate map sends a raw configuration $c$ to
    $(\chi(c|_X),c|_{\mathcal A})$ and hence sends its basis vector to
    the corresponding physical and auxiliary product-basis vector.
    This defines a unitary identification $T$ from the raw
    memory with the physical factors followed by all auxiliary factors.
    At a party with $M_p=\varnothing$, the physical factor is
    one-dimensional.
\end{definition}

\begin{lemma}[Locality and the terminal physical marginal]
    \label{thm:chain_terminal_coordinates}
    \uses{def:chain_terminal_coordinates}
    \lean{
        TNLean.MPS.AreaLaw.ChainSweep.localConfigurationEquiv_restrict,
        TNLean.MPS.AreaLaw.ChainSweep.physicalDensity_terminalRegrouping
    }
    \leanok
    The coordinate map at party $p$ depends precisely on the main
    configuration on $M_p$ and the unchanged auxiliary configuration on
    $\mathcal A_p$. For every terminal vector $\Psi$, with no product
    or normalization assumption, its physical density after applying
    $T$ and discarding all private registers is
    \begin{align}
        \operatorname{Tr}_{\rm private}|T\Psi\rangle\langle T\Psi|
        =U_\chi\bigl(\operatorname{Tr}_{\mathcal A}
        |\Psi\rangle\langle\Psi|\bigr)U_\chi^*,\notag
    \end{align}
    where the expression on the right uses the original main and
    auxiliary sheet coordinates, and $U_\chi$ is their specified main
    block change of coordinates.
\end{lemma}
\begin{proof}\leanok
    \uses{def:chain_terminal_coordinates}
    Restriction to an owner fibre commutes with the main block
    identification and leaves every auxiliary coordinate unchanged.
    In product orthonormal bases, the density transforms by the
    tensor product of the physical and private changes of basis.
    If $V$ is the private change of orthonormal basis, then
    $V^*V=1$ and $\operatorname{Tr}_{\rm private}((U_\chi\otimes V)
        \rho(U_\chi\otimes V)^*)=U_\chi(\operatorname{Tr}_{\mathcal A}\rho)U_\chi^*$.
    Applying this identity to $\rho=|\Psi\rangle\langle\Psi|$ gives the assertion.
\end{proof}

\begin{lemma}[Realization by private operations]
    \label{thm:chain_terminal_private_operations}
    \uses{def:chain_terminal_coordinates}
    \lean{
        TNLean.MPS.AreaLaw.ChainSweep.terminalRegroupingWord,
        TNLean.MPS.AreaLaw.ChainSweep.terminalRegroupingWord_props,
        TNLean.MPS.AreaLaw.ChainSweep.eval_terminalRegroupingWord
    }
    \leanok
    The unitary $T$ is realized by a finite sequence of private
    operations and changes of tensor-factor order. No nonprivate gate
    is introduced, and the equality holds on every input vector.
\end{lemma}
\begin{proof}\leanok
    \uses{def:chain_terminal_coordinates, lem:lattice_site_register_coordinates}
    First place the main owner fibres in the prescribed party order,
    followed by the complete auxiliary list. These are precisely the
    original labels, each occurring once. At each party apply the local
    main-block identification, leaving all other factors fixed.
    An empty main fibre produces its one-dimensional physical factor.
    The resulting map agrees with $T$ on every raw product-basis
    vector and therefore on the whole memory space by linearity.
\end{proof}

\begin{definition}[The terminal reference vector]
    \label{def:chain_terminal_reference}
    \uses{def:chain_terminal_registers, thm:chain_initial_copies,
        def:chain_terminal_coordinates}
    \lean{
        TNLean.MPS.AreaLaw.ChainSweep.auxiliaryReference,
        TNLean.MPS.AreaLaw.ChainSweep.auxiliaryReference_apply,
        TNLean.MPS.AreaLaw.ChainSweep.terminalAuxiliarySiteOrder,
        TNLean.MPS.AreaLaw.ChainSweep.terminalRawOrder,
        TNLean.MPS.AreaLaw.ChainSweep.terminalRawOrder_nodup,
        TNLean.MPS.AreaLaw.ChainSweep.mem_terminalRawOrder,
        TNLean.MPS.AreaLaw.ChainSweep.terminalAuxiliaryFrames,
        TNLean.MPS.AreaLaw.ChainSweep.terminalFrames,
        TNLean.MPS.AreaLaw.ChainSweep.stackRegs_terminalFrames,
        TNLean.MPS.AreaLaw.ChainSweep.terminalStackReadout
    }
    \leanok
    Retain the chain, parties and labelled copies of
    Definition~\ref{def:chain_terminal_registers}, with $q\ge1$.
    Let $\mathcal C$ be the auxiliary copy family, ordered as in
    Lemma~\ref{thm:chain_initial_copies}. For an auxiliary configuration
    $a:\mathcal A\to\{0,\ldots,q-1\}$, put $a_c(x)=a(c,x)$.
    For an arbitrary chain vector $\Omega$, define the auxiliary vector by
    \begin{align}
        \Xi(a)=\prod_{c\in\mathcal C}\Omega(a_c).\notag
    \end{align}
    The empty copy family gives the scalar $\Xi=1$.
    Each copy retains its own site labels and retired owners; a copy
    may be entangled across its owners.

    Let $\mathcal F_{\rm term}$ be the list consisting of the main raw
    sheet with its prescribed main owner, followed by the raw auxiliary
    sheets at their prescribed retired owners in this copy order.
    Write $v_{\rm ref}$ for the product of their raw copy reference
    vectors associated with $\Omega$ in this specified memory. Enumerate the
    main labels first, followed
    by every copy-site label in the prescribed copy order. Each original
    raw label occurs once, and flattening the sheet list gives exactly
    this register list. Denote by $\mathcal R$ the main/auxiliary
    sheet readout transported along that register-list equality.
\end{definition}

\begin{lemma}[The reference in its sheet coordinates]
    \label{thm:chain_terminal_reference_coordinates}
    \uses{def:chain_terminal_reference}
    \lean{TNLean.MPS.AreaLaw.ChainSweep.terminalStackReadout_framesVec}
    \leanok
    For every chain vector $\Omega$,
    \begin{align}
        \mathcal R v_{\rm ref}=\Omega\otimes\Xi.\notag
    \end{align}
\end{lemma}
\begin{proof}\leanok
    \uses{def:chain_terminal_reference, lem:lattice_site_register_coordinates}
    In the raw site basis, a copy has coefficient $\Omega(b)$ at the
    configuration $b$. The tensor product over the labelled copies
    therefore has coefficient $\Omega(b)\prod_{c\in\mathcal C}\Omega(a_c)$
    at $(b,a)$ after flattening. The specified sheet readout sends that
    raw basis vector to the basis vector at $(b,a)$. Its coefficient
    is consequently $\Omega(b)\Xi(a)$, as asserted.
\end{proof}

\begin{lemma}[Normalization and the main reference density]
    \label{thm:chain_terminal_reference_density}
    \uses{def:chain_terminal_reference}
    \lean{
        TNLean.MPS.AreaLaw.ChainSweep.norm_auxiliaryReference,
        TNLean.MPS.AreaLaw.ChainSweep.partialTrace_terminalStackReference
    }
    \leanok
    If $\|\Omega\|=1$, then $\|\Xi\|=1$, including when
    $\mathcal C=\varnothing$, and
    \begin{align}
        \operatorname{Tr}_{\mathcal A}
        |\mathcal Rv_{\rm ref}\rangle\langle\mathcal Rv_{\rm ref}|
        =|\Omega\rangle\langle\Omega|.\notag
    \end{align}
\end{lemma}
\begin{proof}\leanok
    \uses{def:chain_terminal_reference,
        thm:chain_terminal_reference_coordinates}
    The tensor norm gives $\|\Xi\|=\prod_{c\in\mathcal C}\|\Omega\|=1$.
    Tracing the auxiliary factor of the reference in sheet coordinates
    gives $\|\Xi\|^2|\Omega\rangle\langle\Omega|$.
\end{proof}

\begin{lemma}[Density error from a terminal vector error]
    \label{thm:chain_terminal_reference_error}
    \uses{def:chain_terminal_reference}
    \lean{TNLean.MPS.AreaLaw.ChainSweep.rectangularTraceNorm_terminalStackReadout_sub_le}
    \leanok
    Let $\|\Omega\|=1$ and let $\Psi$ belong to the specified terminal
    memory, with $\|\Psi\|\le1$. For every real $\varepsilon$ satisfying
    $\|\Psi-v_{\rm ref}\|\le\varepsilon$, the main density in the
    original sheet coordinates obeys
    \begin{align}
        \left\|\operatorname{Tr}_{\mathcal A}
        |\mathcal R\Psi\rangle\langle\mathcal R\Psi|
        -|\Omega\rangle\langle\Omega|\right\|_1
        \le2\varepsilon.\notag
    \end{align}
\end{lemma}
\begin{proof}\leanok
    \uses{thm:chain_terminal_reference_coordinates,
        thm:chain_terminal_reference_density}
    The readout is an isometry, so
    $\|\mathcal R\Psi-\Omega\otimes\Xi\|
        =\|\Psi-v_{\rm ref}\|\le\varepsilon$ and $\|\mathcal R\Psi\|\le1$.
    The absolute pure-density estimate after partial trace bounds the
    density error by
    $2\|\mathcal R\Psi-\Omega\otimes\Xi\|+1-\|\Xi\|^2$.
    Substituting $\|\Xi\|=1$ gives the result.
\end{proof}

\section{Birth, death and exchange of plain sheets}

\begin{lemma}[Birth and death of a plain sheet]
    \label{thm:chain_plain_birth_death}
    \leanok
    \uses{def:chain_model,def:lattice_sheet,def:lattice_bounded_change}
    Let $L,q\ge1$ be integers, let $p\in\mathbb R$, and let $\Omega$ be a
    unit vector on the chain. A plain sheet has one register $\mathbb C^q$
    per site and an owner assignment $o$. Write $\Omega_o$ for $\Omega$
    placed on those registers. Let $T$ be owned by $P$, put
    $E=\{x:o(x)\ne P\}$, and suppose $I_\Omega(T:E)\le L^{-2p}$.
    For any party $Q$ define
    \begin{align}
        o'(x)=\begin{cases}
                  Q,    & x\in T,    \\
                  o(x), & x\notin T.
              \end{cases}\notag
    \end{align}
    There is a birth from $o$ to $o'$ consisting of private contractions
    and at most one normalized pair source, involving only $P,Q$.
    There is also a death from $o'$ to $o$ using private contractions,
    at most one normalized pair effect and no pair source, again involving
    only $P,Q$. Their maps $\mathcal B,\mathcal D$ satisfy
    \begin{align}
        \|\mathcal B\Omega_o-\Omega_{o'}\| & \le L^{-p},\notag \\
        \|\mathcal D\Omega_{o'}-\Omega_o\| & \le L^{-p}.\notag
    \end{align}
    The information condition for the death is evaluated in its
    after-death assignment $o$. In particular, every fixed $p>0$ gives
    the birth and death estimates of \cite[Part~II]{TNLean2026OneDimensional}.
\end{lemma}
\begin{proof}
    \leanok
    \uses{thm:lattice_sheet_birth_bounded}
    This is Theorem~\ref{thm:lattice_sheet_birth_bounded} on the chain of
    $L$ sites, with the real parameter $L\ge0$ and the exponent $p$. The
    environment of a sheet is exactly $E$. The reversed construction uses
    the information environment after the death.
\end{proof}

\begin{lemma}[Exchange of two plain sheets]
    \label{thm:chain_plain_exchange}
    \leanok
    \uses{def:chain_model,def:lattice_sheet,def:lattice_bounded_change}
    Let $L,q\ge1$ be integers, let $p\in\mathbb R$, and let $\Omega$ be a
    unit vector on the chain. Consider two labelled copies of the sheet with
    owner assignments $o_1,o_2$, a set of sites $Y$, and a party $P$. Put
    \begin{align}
        Z & =\{x:\neg(o_1(x)=P\ \text{and}\ o_2(x)=P)\},\notag \\
        T & =Z\cap Y,\notag                                    \\
        E & =Z\setminus Y.\notag
    \end{align}
    Suppose $I_\Omega(T:E)\le L^{-2p}$. Exchange the two owner assignments
    on $Y$ and leave them unchanged outside $Y$, obtaining $o_1',o_2'$.
    There is a contraction $\mathcal X$ between the corresponding
    register spaces, consisting of owner-preserving reorderings and
    private contractions at $P$, with no shared pair source or effect, such
    that
    \begin{align}
        \bigl\|\mathcal X(\Omega_{o_1}\otimes\Omega_{o_2})
        -\Omega_{o_1'}\otimes\Omega_{o_2'}\bigr\|
        \le4L^{-p}.\notag
    \end{align}
    Every register retains its party; only its sheet label changes in the
    reordering. The two copies remain distinct even when their owner
    assignments agree. For $p>0$, this includes the interval exchange of
    \cite[Part~II]{TNLean2026OneDimensional}.
\end{lemma}
\begin{proof}
    \leanok
    \uses{thm:lattice_sheet_exchange_bounded}
    This is Theorem~\ref{thm:lattice_sheet_exchange_bounded} on the chain of
    $L$ sites, with the real parameter $L\ge0$ and the exponent $p$. Its
    splitting regions are $T,E$. The returned monomial consists of the register
    reordering and the private correction at $P$, and uses no pair resource.
\end{proof}

\section{One-band dissolution}

\begin{definition}[Assignments in one band dissolution]
    \label{def:chain_one_band_assignments}
    \lean{TNLean.MPS.AreaLaw.OneBand.main₀,
        TNLean.MPS.AreaLaw.OneBand.aux₀,
        TNLean.MPS.AreaLaw.OneBand.aux₁,
        TNLean.MPS.AreaLaw.OneBand.aux₂,
        TNLean.MPS.AreaLaw.OneBand.aux₃,
        TNLean.MPS.AreaLaw.OneBand.exchangeRegion,
        TNLean.MPS.AreaLaw.OneBand.exchangeTSet,
        TNLean.MPS.AreaLaw.OneBand.exchangeESet,
        TNLean.MPS.AreaLaw.OneBand.main₄,
        TNLean.MPS.AreaLaw.OneBand.aux₄,
        TNLean.MPS.AreaLaw.OneBand.main₅,
        TNLean.MPS.AreaLaw.OneBand.main₆,
        TNLean.MPS.AreaLaw.OneBand.main₇,
        TNLean.MPS.AreaLaw.OneBand.main₈,
        TNLean.MPS.AreaLaw.OneBand.frames,
        TNLean.MPS.AreaLaw.OneBand.participants}
    \leanok
    \uses{def:lattice_sheet}
    Write $[u,v)_L=\{x\in\{0,\ldots,L-1\}:u\le x<v\}$.
    Fix natural numbers $L,a,w$ with $w\ge1$ and $8w\le a$, a local
    dimension $q\in\mathbb N$, and parties $A,B,C$.
    A sheet is specified by an owner for each site.
    For such a sheet $F$, write $F[T\to P]$ for the assignment obtained
    by giving $T$ to $P$, and write
    $E(F,P)=\{x:\text{the owner of }x\text{ in }F\text{ differs from }P\}$.
    Let $M_0$ have an arbitrary owner function $o$, and let $U_0$ assign
    every site of a distinct auxiliary copy to $C$. Set
    \begin{align}
        U_1 & =U_0[\lbrack a-6w,a-w)_L\to B],        &
        U_2 & =U_1[\lbrack a-5w,a-2w)_L\to A],\notag   \\
        U_3 & =U_2[\lbrack a-4w,a-3w)_L\to C],       &
        Y   & =\{x:a-7w<x<a-\tfrac72w\}.\notag
    \end{align}
    Exchange the ownership assignments of $M_0,U_3$ on $Y$ to define
    $M_4,U_4$, respectively; outside $Y$ retain their assignments.
    The set $Y$ is the sample of an open real interval, also when $w$ is odd.
    Define the remaining main assignments by
    \begin{align}
        M_5 & =M_4[\lbrack a-4w,a)_L\to A],          &
        M_6 & =M_5[\lbrack a-5w,a+w)_L\to B],\notag    \\
        M_7 & =M_6[\lbrack a-3w,a)_L\to C],          &
        M_8 & =M_7[\lbrack a-6w,a-3w)_L\to C].\notag
    \end{align}
    The stage pairs are $(M_0,U_j)$ for $0\le j\le3$, then
    $(M_j,U_4)$ for $4\le j\le8$, and remain $(M_8,U_4)$ thereafter.
    Thus the eight transitions are
    three auxiliary births, one exchange, two main deaths, one main birth,
    and one main death. The auxiliary copy remains unchanged after the exchange.
    Every copy is a distinct occurrence in the list of sheets, including
    copies with the same geometric centre.
    These are the half-open integer bands of
    \cite[Part~II, Construction~\texttt{con:band}]{TNLean2026OneDimensional},
    with the endpoint and retained-register corrections of
    \cite{gap:tnlean_one_dimensional_band_sweep_corrections}.
\end{definition}

\begin{lemma}[Endpoints and information regions of one band]
    \label{thm:chain_one_band_geometry}
    \lean{TNLean.MPS.AreaLaw.OneBand.main₈_eq_changeOwner,
        TNLean.MPS.AreaLaw.OneBand.aux₄_eq_raw,
        TNLean.MPS.AreaLaw.OneBand.exchange_regions}
    \leanok
    \uses{def:chain_one_band_assignments}
    Suppose $o(x)=C$ on $[a-8w,a)_L$.
    The retired auxiliary assignment is
    \begin{align}
        o_{U_4}(x) & =
        \begin{cases}
            A, & x\in[a-3w,a-2w)_L, \\
            B, & x\in[a-2w,a-w)_L,  \\
            C, & \text{otherwise}.
        \end{cases}\notag
    \end{align}
    Moreover $M_8=M_0[\lbrack a,a+w)_L\to B]$, with no assumption on the
    original owners in $\lbrack a,a+2w)_L$.
    Thus the original main ownership is restored outside the initial band;
    this assertion concerns the main sheet, not the retained auxiliary registers.

    For the exchange identities, assume only $w\ge1$, $8w\le a$,
    $o(x)=C$ on $[a-8w,a)_L$, $A\ne C$ and $B\ne C$, and put
    $Z=E(M_0,C)\cup E(U_3,C)$.
    The exchange information regions $T=Z\cap Y$ and $E=Z\setminus Y$ are exactly
    \begin{align}
        T & =[a-6w,a-4w)_L,                   &
        E & =[a-3w,a-w)_L\cup E(M_0,C).\notag
    \end{align}
    The endpoint assertions do not require the three parties to be distinct;
    the two inequalities just stated are required for the exchange-region identities.
\end{lemma}
\begin{proof}\leanok
    \uses{def:chain_one_band_assignments,def:lattice_sheet}
    Before the exchange, the auxiliary owners are $B,A,C,A,B$ on the
    consecutive bands with endpoints $a-6w,a-5w,\ldots,a-w$, and $C$ elsewhere.
    The open window contains $[a-6w,a-4w)_L$ and is contained in
    $[a-7w,a-3w)_L$, where the original main owner is $C$.
    It therefore transfers precisely the two left auxiliary bands to the main
    sheet and restores $C$ on the auxiliary throughout the window. The two
    right auxiliary bands are retained. Substituting the four subsequent
    main assignments gives the main endpoint. For the information regions,
    the sites not held by $C$ on both sheets consist of the two left
    auxiliary bands, the two right auxiliary bands and $E(M_0,C)$.
    The first lie in $Y$; the latter two sets avoid $Y$. This gives $T$ and $E$.
\end{proof}

\begin{lemma}[One-band changes under primitive information bounds]
    \label{thm:chain_one_band_conditional_changes}
    \lean{TNLean.MPS.AreaLaw.OneBand.circuits_of_primitive_information,
        TNLean.MPS.AreaLaw.OneBand.boundary_change_of_primitive_information}
    \leanok
    \uses{def:chain_one_band_assignments,def:lattice_bounded_change,
        def:peps_circuit_source_bound,def:lattice_sheet_stack}
    Assume $o(x)=C$ on $\lbrack a-8w,a)_L$ and $o(x)=A$ on $[a,a+w)_L$; the
    parties $A,B,C$ need not be distinct. Let $q\ge1$ and let $\Omega$ be a unit vector on the
    original chain. Put $\delta=L^{-12}$.
    Assume $I_\Omega(T_j:E_j)\le L^{-24}$ for each row of the following table:
    \begin{align}
        \begin{array}{cccc}
            j & T_j           & E_j      & \text{nonprivate participants} \\
            0 & [a-6w,a-w)_L  & E(U_0,C) & \{C,B\}                        \\
            1 & [a-5w,a-2w)_L & E(U_1,B) & \{B,A\}                        \\
            2 & [a-4w,a-3w)_L & E(U_2,A) & \{A,C\}                        \\
            3 & T             & E        & \{C\}                          \\
            4 & [a-4w,a)_L    & E(M_5,A) & \{C,A\}                        \\
            5 & [a-5w,a+w)_L  & E(M_6,B) & \{A,B\}                        \\
            6 & [a-3w,a)_L    & E(M_6,B) & \{B,C\}                        \\
            7 & [a-6w,a-3w)_L & E(M_8,C) & \{B,C\}
        \end{array}\notag
    \end{align}
    Here $T,E$ are the exact exchange regions above. In particular
    $E(U_0,C)=\varnothing$. All information quantities concern $\Omega$,
    not an evolving approximate vector. The environments in rows $4,5,7$
    are evaluated after the corresponding death; row $6$ uses the sheet
    before the birth.

    Place the chosen adjacent main and auxiliary occurrences between any two
    fixed lists of spectator sheets. Append any further registers with vector
    $\xi$, where $\lVert\xi\rVert\le1$.
    Let $\Psi_j$ be the tensor product, in this list order, of the reference
    vectors of all the sheets in stage $j$, followed by $\xi$; each sheet
    reference is formed from the same original $\Omega$.
    There are eight circuits $c_j$ between these register lists with
    at most one nonprivate gate each, its participants contained in the set
    in row $j$. Each gate expansion has at most one monomial, coefficient
    sum at most one, at most one normalized pair effect and at most one
    normalized pair source per monomial. Their reference errors satisfy
    \begin{align}
        \lVert c_j\Psi_j-\Psi_{j+1}\rVert
         & \le\begin{cases}4\delta,&j=3,\\ \delta,&j\ne3.\end{cases}\notag
    \end{align}
    The exchange circuit $c_3$ has no nonprivate gate. All spectator
    occurrences, the retired auxiliary and the further registers are
    retained at the stated stages and owners.

    At the left boundary there is also a single-death version, for arbitrary
    natural $L,w$ and parties $A,B$. Suppose $M$ is a sheet whose
    owner is $A$ on $T_0=\lbrack 0,w)_L$, and put $M'=M[T_0\to B]$.
    For $q\ge1$, a unit vector $\Omega$, and
    $I_\Omega(T_0:E(M',B))\le L^{-24}$, there is a circuit from $M$ to $M'$
    with nonprivate participants in $\{A,B\}$ and the same bounds one on
    gates, monomials, coefficient sum, pair effects and pair sources, with
    reference error at most $L^{-12}$. It may be placed among the same
    arbitrary spectator sheets and further registers. This boundary assertion
    imposes neither the interior margin hypotheses nor distinctness of $A,B$,
    and introduces no left auxiliary copy.
\end{lemma}
\begin{proof}\leanok
    \uses{thm:chain_plain_birth_death,thm:chain_plain_exchange,
        thm:lattice_bounded_change_closure,
        thm:peps_original_produced_resource_bounds,def:lattice_sheet_stack}
    Apply the birth and death estimates at $p=12$.
    Every modified band is owned by the departing party in its sheet.
    Apply the birth estimate to rows $0,1,2,6$ and its reversed
    death estimate to rows $4,5,7$, with the environments displayed in the
    table. For a death, changing the owner back in its after-death sheet
    recovers the original sheet. The exchange estimate applies to row $3$,
    giving a private circuit with error $4\delta$.
    The register realizations use at most one pair source and at most one
    pair effect per monomial; the exchange uses neither. Placing these circuits beside spectator
    sheets and further registers does not increase their resource bounds or
    their errors: the spectator references and $\xi$ have norm at most one.
    The same death argument gives the boundary version.
\end{proof}

\section{Ownership and labelled copies during the chain sweep}

The assignments follow the sweep construction in
\cite[Part~II]{TNLean2026OneDimensional}, with the interval and
retained-register correction of
\cite{gap:tnlean_one_dimensional_band_sweep_corrections}.
The objects below are owner assignments and ordered lists of sheets.

\begin{definition}[Main-sheet assignments at a stage]
    \label{def:chain_sweep_stage_assignments}
    \lean{
        TNLean.MPS.AreaLaw.ChainSweep.stageOwner,
        TNLean.MPS.AreaLaw.ChainSweep.stageFrame,
        TNLean.MPS.AreaLaw.ChainSweep.remainderBirthTarget,
        TNLean.MPS.AreaLaw.ChainSweep.blockBirthTarget,
        TNLean.MPS.AreaLaw.ChainSweep.leftBandTarget,
        TNLean.MPS.AreaLaw.ChainSweep.rightBandTarget,
        TNLean.MPS.AreaLaw.ChainSweep.afterRemainderBirth,
        TNLean.MPS.AreaLaw.ChainSweep.afterBlockBirth,
        TNLean.MPS.AreaLaw.ChainSweep.afterLeftTransfer,
        TNLean.MPS.AreaLaw.ChainSweep.afterRightTransfer,
        TNLean.MPS.AreaLaw.ChainSweep.leftBandMain,
        TNLean.MPS.AreaLaw.ChainSweep.rightBandMain
    }
    \leanok
    \uses{def:chain_terminal_registers,def:chain_one_band_assignments}
    Let $L,w,q\in\mathbb N$, $X=\{0,\ldots,L-1\}$ and $t=20w$.
    Put $m=\max\{\lfloor L/t\rfloor-3,0\}$ when $w>0$, and $m=0$
    when $w=0$. Use the distinct parties $R_0,\ldots,R_m$ and
    $P_1,\ldots,P_m$ and the labelled auxiliary copies of
    Definition~\ref{def:chain_terminal_registers}.
    For $0\le k\le m$, let $M_k$ be the sheet whose owner is
    $P_{\lfloor x/t\rfloor+1}$ at $x<kt$ and $R_k$ at $x\ge kt$.
    The first case is empty when $t=0$.

    For $0\le k<m$, set $a=kt$, $b=(k+1)t$ and define
    \begin{align}
        T_k^R & =\lbrack b+w,L)_L, & T_k^P & =\lbrack a+w,b)_L,\notag \\
        H_k^- & =\lbrack a,a+w)_L, & H_k^+ & =\lbrack b,b+w)_L.\notag
    \end{align}
    Starting from $M_k$, give $T_k^R$ to $R_{k+1}$, then $T_k^P$ to
    $P_{k+1}$, obtaining sheets $M_k^{(1)},M_k^{(2)}$.
    Define $M_k^{(L)}$ by giving $H_k^-$ in $M_k^{(2)}$ to $P_{k+1}$,
    and define $M_k^{(R)}$ by giving $H_k^+$ in $M_k^{(L)}$ to $R_{k+1}$.

    For $k>0$, also form the main endpoint $\widehat M_k^{(L)}$ of the
    assignments of Definition~\ref{def:chain_one_band_assignments}, with
    initial main sheet $M_k^{(2)}$, centre $a$ and
    $(A,B,C)=(R_k,P_{k+1},P_k)$.
    For $k=0$, instead put $\widehat M_0^{(L)}=M_0^{(L)}$; no party
    $P_0$ or left auxiliary copy is introduced.
    Form $\widehat M_k^{(R)}$ from initial main sheet
    $\widehat M_k^{(L)}$, centre $b$ and
    $(A,B,C)=(R_k,R_{k+1},P_{k+1})$ by the same one-band assignments.
    All integer-ended intervals are half-open, with the correction of
    \cite{gap:tnlean_one_dimensional_band_sweep_corrections}.
\end{definition}

\begin{lemma}[Birth environments, band flanks and main endpoints]
    \label{thm:chain_sweep_stage_owners}
    \lean{
        TNLean.MPS.AreaLaw.ChainSweep.stageOwner_zero,
        TNLean.MPS.AreaLaw.ChainSweep.stageOwner_last,
        TNLean.MPS.AreaLaw.ChainSweep.birth_targets_owned,
        TNLean.MPS.AreaLaw.ChainSweep.birth_environments,
        TNLean.MPS.AreaLaw.ChainSweep.left_band_owners,
        TNLean.MPS.AreaLaw.ChainSweep.first_left_death_environment,
        TNLean.MPS.AreaLaw.ChainSweep.right_band_owners,
        TNLean.MPS.AreaLaw.ChainSweep.afterRightTransfer_eq_stageFrame,
        TNLean.MPS.AreaLaw.ChainSweep.leftBandMain_eq_afterLeftTransfer,
        TNLean.MPS.AreaLaw.ChainSweep.rightBandMain_eq_afterRightTransfer,
        TNLean.MPS.AreaLaw.ChainSweep.rightBandMain_eq_stageFrame
    }
    \leanok
    \uses{def:chain_sweep_stage_assignments}
    The initial sheet assigns the entire main sheet to $R_0$, and $M_m$
    has the terminal main ownership of Definition~\ref{def:chain_terminal_registers}.
    Suppose $w\ge1$ and $0\le k<m$. Each birth target is held by $R_k$
    immediately before that birth. Its environment outside $R_k$ is
    \begin{align}
        E(M_k,R_k)       & =[0,a)_L,                                  &
        E(M_k^{(1)},R_k) & =\lbrack 0,a)_L\cup\lbrack b+w,L)_L.\notag
    \end{align}
    For $k>0$, the owners in $M_k^{(2)}$ on the three intervals
    $\lbrack a-8w,a)_L$, $\lbrack a,a+w)_L$ and $\lbrack a+w,a+2w)_L$ are respectively
    $P_k,R_k,P_{k+1}$.
    In $M_k^{(L)}$, the owners on $\lbrack b-8w,b)_L$, $\lbrack b,b+w)_L$ and
    $[b+w,b+2w)_L$ are respectively $P_{k+1},R_k,R_{k+1}$.
    When $m>0$, the environment at the first stage for the direct death, evaluated
    after that death, is $E(M_0^{(L)},P_1)=[t,L)_L$.
    The four specified owner changes give $M_k^{(R)}=M_{k+1}$.
    If $L\ge3t$, the one-band main endpoints satisfy
    $\widehat M_k^{(L)}=M_k^{(L)}$ and
    $\widehat M_k^{(R)}=M_k^{(R)}=M_{k+1}$.
\end{lemma}
\begin{proof}\leanok
    \uses{def:chain_sweep_stage_assignments,thm:chain_stage_margins,
        thm:chain_one_band_geometry}
    Membership in a completed block fixes the quotient $\lfloor x/t\rfloor$.
    Neither birth changes an earlier completed block. The first target
    lies beyond $b+w$, and the second lies between $a+w$ and $b$;
    substituting their owner assignments gives the two environments.
    The intervals required for the left flank lie in the preceding
    completed block, and those for the right flank lie in the new block.
    The prescribed margins permit the one-band endpoint identity at
    each interior band. At $k=0$ the left endpoint is its direct owner
    change by definition. Checking the completed prefix, the new block
    and the remaining suffix gives $M_k^{(R)}=M_{k+1}$.
\end{proof}

\begin{lemma}[Retired auxiliary owners at their original labels]
    \label{thm:chain_sweep_retired_copies}
    \lean{
        TNLean.MPS.AreaLaw.ChainSweep.rightRetiredAuxiliary,
        TNLean.MPS.AreaLaw.ChainSweep.leftRetiredAuxiliary,
        TNLean.MPS.AreaLaw.ChainSweep.rightRetiredAuxiliary_eq_raw,
        TNLean.MPS.AreaLaw.ChainSweep.leftRetiredAuxiliary_eq_raw
    }
    \leanok
    \uses{def:chain_sweep_stage_assignments,def:chain_terminal_registers}
    Suppose $w\ge1$.
    Denote the right copy at stage $k$ by $\rho_k$, for $0\le k<m$,
    and the left copy used at stage $j+1$ by $\lambda_j$, for
    $0\le j<m-1$. The retired right auxiliary comes from the
    initial main sheet $\widehat M_k^{(L)}$ and has centre $b=(k+1)t$.
    It assigns $[b-3w,b-2w)_L$ to $R_k$,
    $[b-2w,b-w)_L$ to $R_{k+1}$, and the remaining sites to $P_{k+1}$.
    The retired left auxiliary $\lambda_j$ is formed from
    $M_{j+1}^{(2)}$ and has centre $c=(j+1)t$.
    It assigns $[c-3w,c-2w)_L$ to $R_{j+1}$,
    $[c-2w,c-w)_L$ to $P_{j+2}$, and the remaining sites to $P_{j+1}$.
    These are exactly the auxiliary owner functions of
    Definition~\ref{def:chain_terminal_registers}, at the same original
    copy labels. In particular $\rho_j$ and $\lambda_j$ remain distinct
    when their centres coincide.
\end{lemma}
\begin{proof}\leanok
    \uses{thm:chain_sweep_stage_owners,thm:chain_stage_margins,
        thm:chain_one_band_geometry}
    Use the derived flanks for the right band and for the positive-stage
    left band. The one-band auxiliary endpoint retains its two right-hand
    strips. The left margins ensure that natural subtraction in these
    intervals equals integer subtraction, giving the stated owner
    functions without changing either copy label.
\end{proof}

\begin{definition}[The labelled sheet history]
    \label{def:chain_sweep_labelled_history}
    \lean{
        TNLean.MPS.AreaLaw.ChainSweep.auxiliaryStage,
        TNLean.MPS.AreaLaw.ChainSweep.retiredAuxiliaryFrame,
        TNLean.MPS.AreaLaw.ChainSweep.auxiliaryFrameBefore,
        TNLean.MPS.AreaLaw.ChainSweep.sweepFrames,
        TNLean.MPS.AreaLaw.ChainSweep.leftAuxiliaryFrame,
        TNLean.MPS.AreaLaw.ChainSweep.afterLeftFrames,
        TNLean.MPS.AreaLaw.ChainSweep.afterRightFrames,
        TNLean.MPS.AreaLaw.ChainSweep.leftCopyLabel
    }
    \leanok
    \uses{def:chain_sweep_stage_assignments,thm:chain_sweep_retired_copies,
        thm:chain_initial_copies}
    Retain the initial ordered list $\mathcal C$ of all auxiliary copy
    labels. Its order is fixed throughout. The activation stage is
    $s(\rho_j)=j$ and $s(\lambda_j)=j+1$.
    Let $U_c^\infty$ be the retired sheet of copy $c$, and let
    $U_c^0$ be the sheet assigning that copy entirely to its initial
    holder, namely $P_{j+1}$ for either $c=\rho_j$ or $c=\lambda_j$.
    Before stage $k$, set
    \begin{align}
        U_c^{<k}                    & =
        \begin{cases}
            U_c^\infty, & s(c)<k,    \\
            U_c^0,      & s(c)\ge k,
        \end{cases} &
        \mathcal F_k                & =(M_k,(U_c^{<k})_{c\in\mathcal C}).\notag
    \end{align}
    At $k>0$ the selected left label is $\lambda_{k-1}$.
    After the left transfer, replace exactly that occurrence by
    $U_{\lambda_{k-1}}^\infty$ and take main sheet $\widehat M_k^{(L)}$;
    denote the resulting ordered list by $\mathcal F_k^L$.
    At $k=0$ replace no auxiliary occurrence.
    After the right transfer replace exactly $\rho_k$ by $U_{\rho_k}^\infty$
    and take main sheet $\widehat M_k^{(R)}$, obtaining $\mathcal F_k^R$.
\end{definition}

\begin{lemma}[Initial, intermediate and terminal labelled lists]
    \label{thm:chain_sweep_labelled_history}
    \lean{
        TNLean.MPS.AreaLaw.ChainSweep.auxiliaryStage_lt_stageCount,
        TNLean.MPS.AreaLaw.ChainSweep.auxiliaryFrameBefore_zero,
        TNLean.MPS.AreaLaw.ChainSweep.auxiliaryFrameBefore_last,
        TNLean.MPS.AreaLaw.ChainSweep.sweepFrames_zero,
        TNLean.MPS.AreaLaw.ChainSweep.sweepFrames_last,
        TNLean.MPS.AreaLaw.ChainSweep.leftAuxiliaryFrame_right_unused,
        TNLean.MPS.AreaLaw.ChainSweep.auxiliaryFrameBefore_left_unused,
        TNLean.MPS.AreaLaw.ChainSweep.leftAuxiliaryFrame_zero,
        TNLean.MPS.AreaLaw.ChainSweep.leftAuxiliaryFrame_eq_update,
        TNLean.MPS.AreaLaw.ChainSweep.leftAuxiliaryFrame_spectator,
        TNLean.MPS.AreaLaw.ChainSweep.rightAuxiliaryUpdate_eq_next,
        TNLean.MPS.AreaLaw.ChainSweep.rightAuxiliaryUpdate_spectator,
        TNLean.MPS.AreaLaw.ChainSweep.afterRightFrames_eq_sweepFrames,
        TNLean.MPS.AreaLaw.ChainSweep.sweepFrames_select_auxiliary
    }
    \leanok
    \uses{def:chain_sweep_labelled_history,def:chain_terminal_reference}
    For all natural $L,w,q$, $\mathcal F_0$ is the initial main sheet
    followed by the original unused copies, and $\mathcal F_m$ is exactly
    the specified terminal list with all original auxiliary owners.
    Every copy has activation stage less than $m$.
    Before a positive-stage left transfer, its selected copy is unused;
    after the left transfer, the current right copy is still unused.
    At $k=0$ every auxiliary is unchanged by the left transfer.
    At $k>0$ only $\lambda_{k-1}$ changes in that transfer, and only
    $\rho_k$ changes in the subsequent right transfer. All other
    occurrences retain their sheets and their relative order.
    For every $k<m$, one has $\mathcal F_k^R=\mathcal F_{k+1}$.
    Selecting any copy $c$ gives a permutation of $\mathcal F_k$ in which
    the selected sheet and the main sheet are first, followed by the
    other copies in their original order with $c$ omitted.
\end{lemma}
\begin{proof}\leanok
    \uses{def:chain_sweep_labelled_history,thm:chain_sweep_retired_copies,
        thm:chain_sweep_stage_owners}
    Initially no activation stage is negative. Before the last boundary
    every copy's activation stage has been passed, so the retired-sheet
    identities give the terminal auxiliary list. Comparing the two
    activation times $j$ and $j+1$ shows that the left update changes
    precisely $\lambda_{k-1}$ and that the right update changes precisely
    $\rho_k$. The main endpoint identity and these single-label updates
    give the next entire list. For any stage or copy, $m>0$
    implies $w\ge1$ and $L\ge3t$; thus the endpoint comparisons require
    no extra hypothesis here. Removing a selected label from the fixed
    list and placing it first preserves all other occurrences and their order.
\end{proof}

\begin{definition}[Local boundaries and the chronological sheet sequence]
    \label{def:chain_sweep_frame_path}
    \lean{
        TNLean.MPS.AreaLaw.ChainSweep.leftBandFrames,
        TNLean.MPS.AreaLaw.ChainSweep.rightBandFrames,
        TNLean.MPS.AreaLaw.ChainSweep.stageBoundaryFrames,
        TNLean.MPS.AreaLaw.ChainSweep.sweepOperationCount,
        TNLean.MPS.AreaLaw.ChainSweep.operationFrames
    }
    \leanok
    \uses{def:chain_sweep_labelled_history,def:chain_one_band_assignments}
    For a stage $0\le k<m$ with $k>0$, let $\mathcal L_{k,r}$ consist of the $r$th left-band pair
    from Definition~\ref{def:chain_one_band_assignments}, with original
    label $\lambda_{k-1}$, followed by all other copies in their original
    order and their sheets before stage $k$.
    Let $\mathcal R_{k,r}$ be the analogous right-band list with selected
    label $\rho_k$, whose spectators have their sheets after the left
    transfer. In each list the main sheet precedes the selected auxiliary.

    There are nineteen boundaries $\mathcal B_{k,0},\ldots,\mathcal B_{k,18}$
    for eighteen consecutive steps. They begin with $\mathcal F_k$,
    followed by the first and second birth sheets with the auxiliary
    family $U^{<k}$ unchanged, at boundaries $1$ and $2$.
    For $3\le j<10$ and $k>0$, use $\mathcal B_{k,j}=\mathcal L_{k,j-2}$.
    At boundary $10$ use $\mathcal F_k^L$.
    At the first stage use $\mathcal B_{0,j}=\mathcal F_0^L$ for
    $3\le j\le10$: boundary $3$ is the direct death and the following
    seven steps leave the sheet list unchanged.
    For $11\le j<18$ use $\mathcal B_{k,j}=\mathcal R_{k,j-10}$.
    At and beyond boundary $18$ use $\mathcal F_k^R$.

    Put $N=18m$. For a natural index $n$, define
    \begin{align}
        \mathcal G_n & =
        \begin{cases}
            \mathcal B_{\lfloor n/18\rfloor,\,n\bmod18}, & n<N,    \\
            \mathcal F_m,                                & n\ge N.
        \end{cases}\notag
    \end{align}
\end{definition}

\begin{lemma}[Band entrances, exits and the exceptional first stage]
    \label{thm:chain_sweep_local_boundaries}
    \lean{
        TNLean.MPS.AreaLaw.ChainSweep.stageBoundaryFrames_zero,
        TNLean.MPS.AreaLaw.ChainSweep.stageBoundaryFrames_one,
        TNLean.MPS.AreaLaw.ChainSweep.stageBoundaryFrames_two,
        TNLean.MPS.AreaLaw.ChainSweep.stageBoundaryFrames_left,
        TNLean.MPS.AreaLaw.ChainSweep.stageBoundaryFrames_first_left,
        TNLean.MPS.AreaLaw.ChainSweep.stageBoundaryFrames_ten,
        TNLean.MPS.AreaLaw.ChainSweep.stageBoundaryFrames_first_death,
        TNLean.MPS.AreaLaw.ChainSweep.stageBoundaryFrames_right,
        TNLean.MPS.AreaLaw.ChainSweep.stageBoundaryFrames_last,
        TNLean.MPS.AreaLaw.ChainSweep.stageBoundaryFrames_eighteen,
        TNLean.MPS.AreaLaw.ChainSweep.stageBoundaryFrames_left_entrance_perm,
        TNLean.MPS.AreaLaw.ChainSweep.leftBandFrames_eight_perm,
        TNLean.MPS.AreaLaw.ChainSweep.stageBoundaryFrames_right_entrance_perm,
        TNLean.MPS.AreaLaw.ChainSweep.rightBandFrames_eight_perm
    }
    \leanok
    \uses{def:chain_sweep_frame_path}
    The first three boundaries have respectively main sheets
    $M_k,M_k^{(1)},M_k^{(2)}$, with their original ordered auxiliary
    family. At a positive stage, $\mathcal B_{k,2}$ is a permutation of
    $\mathcal L_{k,0}$, and $\mathcal L_{k,8}$ is a permutation of
    $\mathcal F_k^L=\mathcal B_{k,10}$.
    Likewise $\mathcal B_{k,10}$ is a permutation of $\mathcal R_{k,0}$,
    and $\mathcal R_{k,8}$ is a permutation of
    $\mathcal F_k^R=\mathcal B_{k,18}$.
    These permutations select the original copy label and preserve the
    order of all spectator labels; no distinctness of sheet values is required.
    At $k=0$ the transition from boundary $2$ to boundary $3$ changes
    the main owner on $[0,w)_L$ from $R_0$ to $P_1$ and changes no
    auxiliary. Boundaries $3,\ldots,10$ are equal.
    For every $k<m$, $\mathcal B_{k,18}=\mathcal F_{k+1}$, and
    $\mathcal B_{k,j}=\mathcal F_k^R$ for every $j\ge18$.
\end{lemma}
\begin{proof}\leanok
    \uses{def:chain_sweep_frame_path,thm:chain_sweep_labelled_history,
        thm:chain_sweep_stage_owners}
    The interior boundaries are the specified one-band sheets.
    Before its activation the selected label still carries its initial
    constant-owner sheet. Selecting that label therefore gives the
    correct band entrance. At its exit, replacing that one label by
    its retired sheet and restoring the original order gives
    the stated permutation. For the first left transfer the defining
    direct death replaces the entire eight-step band, with seven
    repeated boundaries afterwards. The next-stage identity identifies
    the eighteenth boundary.
\end{proof}

\begin{lemma}[Division into stages and the stationary terminal tail]
    \label{thm:chain_sweep_frame_path}
    \lean{
        TNLean.MPS.AreaLaw.ChainSweep.operationFrames_of_le,
        TNLean.MPS.AreaLaw.ChainSweep.operationFrames_stage,
        TNLean.MPS.AreaLaw.ChainSweep.operationFrames_zero,
        TNLean.MPS.AreaLaw.ChainSweep.operationFrames_last,
        TNLean.MPS.AreaLaw.ChainSweep.operationFrames_stageBoundary,
        TNLean.MPS.AreaLaw.ChainSweep.operationFrames_step
    }
    \leanok
    \uses{def:chain_sweep_frame_path}
    For every $0\le k\le m$, $\mathcal G_{18k}=\mathcal F_k$.
    For every $k<m$ and $0\le j\le18$,
    \begin{align}
        \mathcal G_{18k+j} & =\mathcal B_{k,j}.\notag
    \end{align}
    In particular the two consecutive lists at step $j<18$ are
    $\mathcal B_{k,j}$ and $\mathcal B_{k,j+1}$, including the last
    step of every stage. The sequence starts at the original initial
    list and equals the terminal list for every $n\ge18m$.
    If $m=0$, the initial and terminal lists coincide and the entire
    sequence is constant. If $m=1$, there is one eighteen-step stage,
    with its direct first left death, seven unchanged-sheet steps and
    one right auxiliary copy, followed by the stationary terminal tail.
\end{lemma}
\begin{proof}\leanok
    \uses{def:chain_sweep_frame_path,thm:chain_sweep_local_boundaries,
        thm:chain_sweep_labelled_history}
    For $j<18$, division of $18k+j$ by $18$ gives quotient $k$ and
    remainder $j$. At $j=18$, the eighteenth local boundary is
    $\mathcal F_{k+1}$; the next quotient and zero remainder give that
    same list, or the terminal branch if $k+1=m$.
    This also identifies every complete stage boundary. Beyond $18m$
    the sequence is constant by definition. If $m=0$, the initial and
    terminal identities apply to the same boundary.
\end{proof}

\section{Participants and conditional circuits of the chain stages}

Use the parties, labelled copies and ordered lists of
Definitions~\ref{def:chain_terminal_registers},
\ref{def:chain_sweep_labelled_history} and~\ref{def:chain_sweep_frame_path}.
Thus $t=20w$, there are $m$ stages, and stage $k$ has boundaries
$\mathcal B_{k,0},\ldots,\mathcal B_{k,18}$.
These are the prescribed stages of
\cite[Part~II, the sweep construction]{TNLean2026OneDimensional},
with the retained-copy correction of
\cite{gap:tnlean_one_dimensional_band_sweep_corrections}.

\begin{definition}[Prescribed participants and position counts]
    \label{def:chain_stage_participants}
    \lean{TNLean.MPS.AreaLaw.ChainSweep.stageParticipants,
        TNLean.MPS.AreaLaw.ChainSweep.adjacentStageIndices,
        TNLean.MPS.AreaLaw.ChainSweep.stageOccurrences,
        TNLean.MPS.AreaLaw.ChainSweep.participatingStageIndices,
        TNLean.MPS.AreaLaw.ChainSweep.scheduledParticipationCount}
    \leanok
    \uses{def:chain_sweep_frame_path,def:chain_one_band_assignments}
    Write $\Pi_j(A,B,C)$ for the participant set in row $j$ of the table
    in Lemma~\ref{thm:chain_one_band_conditional_changes}.
    For $0\le k<m$ and $0\le j<18$, prescribe
    \begin{align}
        \mathcal S_{k,j} & =
        \begin{cases}
            \{R_k,R_{k+1}\},                 & j=0,             \\
            \{R_k,P_{k+1}\},                 & j=1,             \\
            \Pi_{j-2}(R_k,P_{k+1},P_k),      & k>0,\ 2\le j<10, \\
            \{P_1,R_0\},                     & k=0,\ j=2,       \\
            \varnothing,                     & k=0,\ 3\le j<10, \\
            \Pi_{j-10}(R_k,R_{k+1},P_{k+1}), & 10\le j<18.
        \end{cases}
        \label{eq:chain_stage_participants}
    \end{align}
    Let $\mathrm{Occ}(p)=\{(k,j):p\in\mathcal S_{k,j}\}$,
    let $J(p)$ be its set of stage indices, and set
    $N_{\rm pos}(p)=|\mathrm{Occ}(p)|$.
    Define $\mathcal A(R_r)=\{\max\{r-1,0\},r\}$ and
    $\mathcal A(P_{r+1})=\{r,r+1\}$.
    The number $N_{\rm pos}(p)$ counts prescribed positions, including
    positions whose realized operation may be private.
\end{definition}

\begin{lemma}[Incidence in two adjacent stages]
    \label{thm:chain_stage_incidence}
    \lean{TNLean.MPS.AreaLaw.ChainSweep.remainder_participating_stage,
        TNLean.MPS.AreaLaw.ChainSweep.completed_participating_stage,
        TNLean.MPS.AreaLaw.ChainSweep.stageIndex_mem_adjacent,
        TNLean.MPS.AreaLaw.ChainSweep.participatingStageIndices_subset,
        TNLean.MPS.AreaLaw.ChainSweep.participatingStageIndices_card_le_two,
        TNLean.MPS.AreaLaw.ChainSweep.scheduledParticipationCount_eq_sum,
        TNLean.MPS.AreaLaw.ChainSweep.scheduledParticipationCount_le_thirtySix,
        TNLean.MPS.AreaLaw.ChainSweep.scheduledParticipationCount_le_forty}
    \leanok
    \uses{def:chain_stage_participants}
    Membership in the prescribed sets implies
    \begin{align}
        R_r\in\mathcal S_{k,j}     & \ \Longrightarrow\ k=r\ \text{or}\ k+1=r,
        \label{eq:chain_stage_remainder_incidence}                             \\
        P_{r+1}\in\mathcal S_{k,j} & \ \Longrightarrow\ k=r\ \text{or}\ k=r+1.
        \label{eq:chain_stage_completed_incidence}
    \end{align}
    Consequently $J(p)\subseteq\mathcal A(p)$ and $|J(p)|\le2$. The position count is
    \begin{align}
        N_{\rm pos}(p)
         & =\sum_{k<m}|\{j<18:p\in\mathcal S_{k,j}\}|\le36\le40.
        \label{eq:chain_stage_position_bound}
    \end{align}
    This is the position count in the sweep construction of
    \cite[Part~II]{TNLean2026OneDimensional}.
\end{lemma}
\begin{proof}\leanok
    \uses{def:chain_stage_participants}
    Every row of a band uses only its three specified parties. Thus the
    possible stage roles are $R_k,R_{k+1},P_{k+1}$ and, only for $k>0$,
    $P_k$. Comparing these labels gives
    (\ref{eq:chain_stage_remainder_incidence}) and
    (\ref{eq:chain_stage_completed_incidence}). Projecting
    $\mathrm{Occ}(p)$ to its stage coordinate gives $J(p)$.
    Each fibre has at most eighteen elements, so
    $|\mathrm{Occ}(p)|\le18|J(p)|\le36$; summing these fibres gives
    (\ref{eq:chain_stage_position_bound}). No predecessor is formed at stage zero.
\end{proof}

\begin{lemma}[Local stage circuits under primitive information bounds]
    \label{thm:chain_conditional_local_circuits}
    \lean{TNLean.MPS.AreaLaw.ChainSweep.local_stage_circuits_of_primitive_information}
    \leanok
    \uses{def:chain_stage_participants,def:chain_sweep_frame_path,
        def:lattice_bounded_change,def:peps_circuit_source_bound,def:lattice_sheet_stack}
    Suppose $q\ge1$, $w\ge1$ and $0\le k<m$.
    Let $\Omega$ be a unit vector on the original chain and put
    $\delta=L^{-12}$. Assume the two before-birth bounds
    \begin{align}
        I_\Omega(T_k^R:E(M_k,R_k))       & \le L^{-24},
        \label{eq:chain_stage_remainder_information}    \\
        I_\Omega(T_k^P:E(M_k^{(1)},R_k)) & \le L^{-24}.
        \label{eq:chain_stage_block_information}
    \end{align}
    At $k=0$, assume also the after-death bound
    \begin{align}
        I_\Omega([0,w)_L:E(M_0^{(L)},P_1)) & \le L^{-24}.
        \label{eq:chain_stage_boundary_information}
    \end{align}
    For $k>0$, assume each of the eight information inequalities in the
    table of Lemma~\ref{thm:chain_one_band_conditional_changes}, with
    \begin{align}
        (M_0,a,A,B,C) & =(M_k^{(2)},kt,R_k,P_{k+1},P_k).
        \label{eq:chain_stage_left_information}
    \end{align}
    For every $k$, assume the eight inequalities in that table with
    \begin{align}
        (M_0,a,A,B,C) & =(\widehat M_k^{(L)},(k+1)t,R_k,R_{k+1},P_{k+1}).
        \label{eq:chain_stage_right_information}
    \end{align}
    In each table the auxiliary starts entirely at its specified party $C$.
    All nineteen premises concern the original $\Omega$: the boundary
    premise is conditional on $k=0$, and the eight left premises on $k>0$.
    In particular rows $4,5,7$ use their after-death environments,
    and row $6$ uses its before-birth environment.

    Let $\psi(\mathcal F)$ denote the tensor product of the sheet references
    formed from $\Omega$ in the order of the list $\mathcal F$.
    There exist circuits
    $c_{k,j}:\mathcal H_{\mathcal B_{k,j}}\to
        \mathcal H_{\mathcal B_{k,j+1}}$, for $0\le j<18$, each with at most
    one nonprivate gate, whose participants lie in $\mathcal S_{k,j}$, and
    \begin{align}
        \|c_{k,j}\psi(\mathcal B_{k,j})-
        \psi(\mathcal B_{k,j+1})\| & \le4\delta.
        \label{eq:chain_stage_local_error}
    \end{align}
    Each original nonprivate gate has at most one recorded monomial,
    absolute coefficient sum at most one, and at most one normalized
    pair source and one normalized pair effect per monomial.
    At stage zero, position $2$ joins the boundaries of the direct
    death into $P_1$; positions $3,\ldots,9$ are identity circuits, with
    zero error and no nonprivate gate. Every original auxiliary label and the order
    of all spectator occurrences are retained.
\end{lemma}
\begin{proof}\leanok
    \uses{thm:chain_sweep_stage_owners,thm:chain_stage_margins,
        thm:chain_one_band_conditional_changes,thm:chain_sweep_local_boundaries,
        thm:chain_plain_birth_death,thm:lattice_bounded_change_closure}
    The two birth targets are held by $R_k$ in their before-birth
    sheets. Apply the birth estimate with
    (\ref{eq:chain_stage_remainder_information}) and
    (\ref{eq:chain_stage_block_information}). The derived band flanks and
    margins permit the two substitutions
    (\ref{eq:chain_stage_left_information}) and
    (\ref{eq:chain_stage_right_information}) in the one-band theorem.
    Select the original copy $\lambda_{k-1}$ only for $k>0$, and select
    $\rho_k$ on the right. The entrance and exit permutations identify
    these band circuits with the prescribed boundaries, restoring the
    original spectator order within the same operation position.
    At $k=0$, the boundary death uses
    (\ref{eq:chain_stage_boundary_information}); the seven subsequent
    boundaries agree. Tensoring with spectator references preserves the
    bounds, since their norms are at most one. Widening each ordinary
    error $\delta$ to $4\delta$ gives (\ref{eq:chain_stage_local_error}).
\end{proof}

\begin{definition}[The conditional chronological middle circuit]
    \label{def:chain_conditional_middle_circuit}
    \lean{TNLean.MPS.AreaLaw.ChainSweep.sweepMiddleStep,
        TNLean.MPS.AreaLaw.ChainSweep.sweepMiddleCircuit}
    \leanok
    \uses{def:chain_sweep_frame_path,def:chain_stage_participants,
        def:lattice_bounded_change,def:peps_circuit_source_bound,
        def:peps_original_chronology_composite}
    Let $q\ge1$ and let $\Omega$ be any original chain vector.
    Write $\delta_L=L^{-12}$ for $L>0$ and $\delta_0=0$.
    Suppose each given circuit $c_{k,j}$ between the successive
    boundary lists is a bounded change from $\psi(\mathcal B_{k,j})$ to
    $\psi(\mathcal B_{k,j+1})$, with participants in $\mathcal S_{k,j}$,
    bounds $r=K=1$ and error $4\delta_L$, and has pair-source bound one.
    At $k=0$ and $3\le j<10$, require the identity circuit, zero
    reference error and zero nonprivate count, as in the local conclusion.
    Put $N=18m$. On the chronological lists $\mathcal G_n$, define
    \begin{align}
        d_n & =
        \begin{cases}
            c_{\lfloor n/18\rfloor,n\bmod18}, & n<N,    \\
            1_{\mathcal H_{\mathcal F_m}},    & n\ge N,
        \end{cases}
            & C & =d_{N-1}\cdots d_0.
        \label{eq:chain_conditional_middle}
    \end{align}
    The boundary identities of Lemma~\ref{thm:chain_sweep_frame_path}
    identify the displayed source and target spaces. For $N=0$, take
    $C$ to be the identity. It maps the memory of the initial list
    $\mathcal F_0$ to that of the exact terminal list $\mathcal F_m$.
\end{definition}

\begin{lemma}[Counts, expansions and error of the middle circuit]
    \label{thm:chain_conditional_middle_bounds}
    \lean{TNLean.MPS.AreaLaw.ChainSweep.sweepMiddleStep_tail_identity,
        TNLean.MPS.AreaLaw.ChainSweep.sweepMiddleStep_first_padding,
        TNLean.MPS.AreaLaw.ChainSweep.sweepMiddleCircuit_participation_le_forty,
        TNLean.MPS.AreaLaw.ChainSweep.sweepMiddleCircuit_nonprivateCount_le,
        TNLean.MPS.AreaLaw.ChainSweep.sweepMiddleCircuit_resource_bounds,
        TNLean.MPS.AreaLaw.ChainSweep.sweepMiddleCircuit_norm_and_error}
    \leanok
    \uses{def:chain_conditional_middle_circuit}
    In Definition~\ref{def:chain_conditional_middle_circuit}, let $N(C)$
    count nonprivate occurrences and let $n_p(C)$ count those
    involving party $p$. Then
    \begin{align}
        N(C) & \le18m, & n_p(C) & \le40.
        \label{eq:chain_middle_circuit_counts}
    \end{align}
    Every original nonprivate gate retains the recorded monomial,
    coefficient-sum, pair-effect and pair-source bounds one.
    If $\|\Omega\|=1$, then, with $\Psi=C\psi(\mathcal F_0)$,
    \begin{align}
        \|\Psi\|                    & \le1,              &
        \|\Psi-\psi(\mathcal F_m)\| & \le4(18m)\delta_L.
        \label{eq:chain_middle_circuit_error}
    \end{align}
    The tail $d_n$ for $n\ge18m$ is the identity. When a first stage
    exists, its positions $3,\ldots,9$ are identities. If $m=0$, the
    middle is the identity on the common initial and terminal list.
    If $m=1$, it has one eighteen-position stage with the direct
    left-death boundaries, seven identities and one right auxiliary copy.
\end{lemma}
\begin{proof}\leanok
    \uses{thm:chain_stage_incidence,thm:peps_original_chronology_composition,
        lem:peps_protocol_telescoping,thm:chain_initial_copies}
    The count identities for chronological composition sum the
    position counts. There is at most one nonprivate occurrence per position,
    and an occurrence involving $p$ lies in a position counted by
    $N_{\rm pos}(p)$. Thus Lemma~\ref{thm:chain_stage_incidence} gives
    (\ref{eq:chain_middle_circuit_counts}). The recorded expansion at
    each occurrence is unchanged by composition or by the boundary
    identifications, preserving its four bounds one.
    Put $x_0=\psi(\mathcal G_0)$ and $x_{n+1}=d_nx_n$.
    The circuits are contractions, so $\|x_N\|\le\|x_0\|=1$ and
    \begin{align}
        \|x_N-\psi(\mathcal G_N)\|
         & \le\sum_{n<N}\|d_n\psi(\mathcal G_n)-\psi(\mathcal G_{n+1})\|
        \le4N\delta_L.\notag
    \end{align}
    The initial and terminal boundary identities give
    (\ref{eq:chain_middle_circuit_error}). The remaining assertions use
    the specified identity slots and stationary terminal lists.
\end{proof}

\begin{definition}[Party placement and tagged prefix crossings]
    \label{def:chain_party_placement}
    \lean{TNLean.MPS.AreaLaw.ChainSweep.partyLevel,
        TNLean.MPS.AreaLaw.ChainSweep.prefixCrossingLinks}
    \leanok
    \uses{def:chain_terminal_registers,def:compression_operator_links,
        def:compression_operator_endpoints}
    Place $R_i$ at level $\ell(R_i)=i$ and $P_{j+1}$ at
    $\ell(P_{j+1})=j$.
    Retain all parties, including $R_m$ and the parties with scalar
    main output. For a tagged star or sample link $e$, write
    $s(e),t(e)$ for its two oriented party endpoints. Let $\partial_r$
    be the links with exactly one endpoint in $\{p:\ell(p)<r\}$.
    Every crossing link retains its original gate and kind tag.
\end{definition}

\begin{lemma}[Derived stage support, arity and placement distance]
    \label{thm:chain_middle_gate_support}
    \lean{TNLean.MPS.AreaLaw.ChainSweep.stageParticipants_level_dist_le_two,
        TNLean.MPS.AreaLaw.ChainSweep.card_le_eight_of_subset_stageParticipants,
        TNLean.MPS.AreaLaw.ChainSweep.sweepMiddleCircuit_participants_subset_stage,
        TNLean.MPS.AreaLaw.ChainSweep.sweepMiddleCircuit_participants_level_dist_le_two,
        TNLean.MPS.AreaLaw.ChainSweep.sweepMiddleCircuit_participants_card_le_forty}
    \leanok
    \uses{def:chain_party_placement,def:chain_conditional_middle_circuit}
    Write $\operatorname{part}(g)$ for the participant set of a gate.
    Every pair of parties in one $\mathcal S_{k,j}$ has level distance
    at most two, and each finite subset of $\mathcal S_{k,j}$ has at
    most eight elements. Each nonprivate gate $g$ of the middle
    circuit has a stage and position with
    \begin{align}
        \operatorname{part}(g) & \subseteq\mathcal S_{k,j}.
        \label{eq:chain_middle_gate_support}
    \end{align}
    Consequently $|\operatorname{part}(g)|\le40$, and
    $|\ell(p)-\ell(q)|\le2$ whenever $p,q\in\operatorname{part}(g)$.
\end{lemma}
\begin{proof}\leanok
    \uses{thm:chain_stage_incidence,thm:peps_original_chronology_composition}
    The explicit roles at stage $k$ have levels between
    $\max\{k-1,0\}$ and $k+1$, giving the distance bound.
    They lie in a window of at most four levels,
    $[\max\{k-1,0\},k+3)$, with at most one remainder party and one
    completed party at each level, giving the coarse bound eight.
    Each local gate has the required stage support. Transporting its
    participant set through the boundary identifications and
    chronological composition gives (\ref{eq:chain_middle_gate_support}).
    Apply the two stage bounds to this same support witness, and use
    $8\le40$. When $m=0$, there is no nonprivate gate to which a stage
    need be assigned.
\end{proof}

\begin{lemma}[Congestion of the tagged links]
    \label{thm:chain_tagged_link_congestion}
    \lean{TNLean.MPS.AreaLaw.ChainSweep.prefixCrossingLinks_zero,
        TNLean.MPS.AreaLaw.ChainSweep.prefixCrossingLinks_eq_empty_of_stageCount_lt,
        TNLean.MPS.AreaLaw.ChainSweep.prefixCrossingLinks_card_le_of_originalCircuit,
        TNLean.MPS.AreaLaw.ChainSweep.compressedMiddle_prefixCrossingLinks_card_le,
        TNLean.MPS.AreaLaw.ChainSweep.compressedMiddle_prefixCrossingLinks_card_le_24960}
    \leanok
    \uses{def:chain_party_placement}
    Let $W$ be an original circuit on these parties such that any two
    participants of one nonprivate gate have level distance at most two.
    Consider a tagged distributed link family on these parties. Suppose each
    party is incident to at most $d$ links and that the two endpoints
    of every link participate in a common original nonprivate gate of
    $W$. Then, for every natural prefix index $r$,
    \begin{align}
        |\partial_r| & \le8d.
        \label{eq:chain_tagged_link_congestion}
    \end{align}
    The endpoint prefixes $r=0$ and $r>m$ have no crossing links.
    For the middle circuit $C$, its derived support supplies the
    required gate distance. Thus the common-original-gate and
    degree conclusions of distributed compression at $d=3120$ imply
    \begin{align}
        |\partial_r| & \le24960.
        \label{eq:chain_middle_link_congestion}
    \end{align}
    Here $3120=2\cdot40\cdot39$ is the degree bound in
    Theorem~\ref{thm:peps_distributed_construction_compression} at $b=40$,
    from \cite[Theorem~5.2]{OpenAI2026PolynomialPEPS}.
    This gives a coarse bound for the routing argument of
    \cite[Part~II, proof of Theorem~B]{TNLean2026OneDimensional}.
\end{lemma}
\begin{proof}\leanok
    \uses{thm:chain_middle_gate_support,def:compression_operator_endpoints}
    A crossing link with endpoint levels differing by at most two has
    both endpoints in
    $\mathcal N_r=\{p:\max\{r-2,0\}\le\ell(p)<r+2\}$.
    Mapping a party to its level offset and its remainder/completed
    type injects $\mathcal N_r$ into four levels times two types, so
    $|\mathcal N_r|\le8$. Every crossing link belongs to the incident
    link set of its first endpoint. The union bound therefore gives
    \begin{align}
        |\partial_r| & \le\sum_{p\in\mathcal N_r}
        |\{e:p\text{ is an endpoint of }e\}|\le8d.\notag
    \end{align}
    This counts the full link tags; parallel links stay distinct and
    a loop never crosses a prefix. Every party level is at most $m$,
    proving the endpoint assertions. For $C$, apply
    Lemma~\ref{thm:chain_middle_gate_support} to the common original
    gate and substitute $d=3120$.
\end{proof}

\section{Conditional preparation and ranks in the original chain coordinates}

Retain the parties, original copy labels, stage boundaries and placement of
Definitions~\ref{def:chain_terminal_registers},
\ref{def:chain_sweep_frame_path}, \ref{def:chain_stage_participants}
and~\ref{def:chain_party_placement}.
For an original circuit $W$, let $N(W)$ count its nonprivate occurrences
and $n_a(W)$ those involving party $a$; write $\operatorname{part}(g)$
for the participant set of a gate.
Write $F_S(v)$ for the coefficient matrix of a party tensor $v$ across
$S\mid S^c$, and $M_k(v)$ for the original-site cut matrix of
Definition~\ref{def:chain_open_mps}.

\begin{lemma}[Preparation from the conditional local family]
    \label{thm:chain_conditional_preparation}
    \lean{TNLean.MPS.AreaLaw.ChainSweep.exists_preparation_of_local_circuits}
    \leanok
    \uses{def:chain_conditional_middle_circuit,def:chain_terminal_reference,
        def:chain_terminal_coordinates}
    Let $L,w,q$ be natural numbers with $L>0$ and $q\ge1$, and let
    $\Omega$ be a unit vector on the original chain. Suppose the circuits $c_{k,j}$ satisfy exactly
    the local conditions in Definition~\ref{def:chain_conditional_middle_circuit}:
    they are bounded changes between the prescribed boundary references,
    with participants in $\mathcal S_{k,j}$, bounds $r=K=1$, error
    $4L^{-12}$ and pair-source bound one; the seven first-stage padding
    positions are identities with zero error and zero nonprivate count.
    Put $N=18m$ and $\varepsilon=4NL^{-12}$.

    Let $\mathcal F_0$ be the original ordered initial sheet list.
    Its sheets have constant owner; let $h(F)$ be the owner of site zero
    in sheet $F$. For each party $a$, define $H_a$ as the tensor product,
    in the original order, of the sheet spaces with $h(F)=a$, and let
    $x_a$ be the product of their $\Omega$-reference vectors. Then
    $\|x_a\|=1$, including an empty product. From the constructed product
    input $X=\bigotimes_a x_a$ there is an original circuit $C$ to the
    terminal physical factors followed by all retained private
    registers. The physical factor at $a$ has dimension $d_a=q^{|M_a|}$.
    Its output $\Psi=CX$ satisfies
    \begin{align}
        N(C) & \le N, & n_a(C) & \le40, & \|\Psi\| & \le1.
        \label{eq:chain_preparation_counts}
    \end{align}
    Each original nonprivate gate has at most one recorded monomial, absolute
    coefficient sum at most one, and at most one normalized pair source
    and one normalized pair effect per monomial. Every nonprivate gate
    $g$ of this same $C$ has a stage and position with
    \begin{align}
        \operatorname{part}(g) & \subseteq\mathcal S_{k,j}.
        \label{eq:chain_preparation_support}
    \end{align}
    In particular its arity is at most $40$ and its participant levels
    have pairwise distance at most two.

    Let $J$ be the prescribed identification of the terminal sheet stack
    with its raw register list, and let $T$ be the private
    terminal regrouping. Put $Y=J^{-1}T^{-1}\Psi$.
    With the sheet readout $\mathcal R$ and copy-labelled vector $\Xi$
    of Definition~\ref{def:chain_terminal_reference}, one has
    \begin{align}
        \|\Xi\|                          & =1,             &
        \|\mathcal RY-\Omega\otimes\Xi\| & \le\varepsilon.
        \label{eq:chain_preparation_sheet_error}
    \end{align}
    The prescribed owner-grouping isometry $A_{\rm aux}$ identifies the
    private memory with the Euclidean space of the original auxiliary
    copy-site configurations, so that memory is finite-dimensional.
    Write $V$ for the readout of $\Psi$ in the prescribed physical
    coordinates and the fixed standard private orthonormal basis, and put
    $\rho=\operatorname{Tr}_{\rm private}|V\rangle\langle V|$.
    For the fixed main-block coordinate map $U_\chi$,
    \begin{align}
        \|\rho-|U_\chi\Omega\rangle\langle U_\chi\Omega|\|_1
         & \le2\varepsilon.
        \label{eq:chain_preparation_physical_error}
    \end{align}
    All auxiliary copy-site labels remain at their specified private
    owners. The vector $\Xi$ is the product over original copies,
    including the scalar $1$ for an empty copy family; it need not
    factor over its final parties. This is the conditional preparation
    step in \cite[Part~II]{TNLean2026OneDimensional}, with
    \cite{gap:tnlean_one_dimensional_band_sweep_corrections}.
\end{lemma}
\begin{proof}\leanok
    \uses{thm:chain_initial_copies,thm:chain_conditional_middle_bounds,
        thm:chain_middle_gate_support,thm:chain_terminal_private_operations,
        thm:chain_terminal_reference_density,thm:chain_terminal_reference_error,
        thm:chain_terminal_coordinates}
    Compose the constructed private initialization, the given family's
    chronological middle, and the private terminal regrouping.
    The two ends have no nonprivate occurrence; each such occurrence
    of the composite is the same occurrence of the middle with the same
    participant set. Its support and bounds therefore give
    (\ref{eq:chain_preparation_counts}) and
    (\ref{eq:chain_preparation_support}). The initialization is exact
    and the terminal regrouping is an isometry, so pulling back by
    $J^{-1}T^{-1}$ gives the middle output. The terminal reference
    identity and its derived auxiliary norm give
    (\ref{eq:chain_preparation_sheet_error}). The private layout
    is finite-dimensional through the grouping by auxiliary
    owner. The physical marginal theorem is invariant under the choice
    of private orthonormal basis, so the terminal density-error estimate
    gives (\ref{eq:chain_preparation_physical_error}) in the stated
    physical coordinates.
\end{proof}

\begin{lemma}[Three-party bound at a stage position]
    \label{thm:chain_stage_arity_three}
    \lean{TNLean.MPS.AreaLaw.ChainSweep.card_le_three_of_subset_stageParticipants}
    \leanok
    \uses{def:chain_stage_participants}
    For arbitrary natural $L,w$, any finite party set $A$ contained in
    one prescribed $\mathcal S_{k,j}$ satisfies $|A|\le3$.
    Consequently every nonprivate gate of the circuit constructed in
    Lemma~\ref{thm:chain_conditional_preparation} has at most three parties.
\end{lemma}
\begin{proof}\leanok
    \uses{def:chain_stage_participants,thm:chain_conditional_preparation}
    The two births and first boundary death have at most two parties,
    the padding sets are empty, and every band row uses a subset of
    its three designated roles. Coincident roles decrease the count.
    Applying this finite-set bound to
    (\ref{eq:chain_preparation_support}) gives the last assertion.
\end{proof}

\begin{lemma}[One useful column and all its cut ranks]
    \label{thm:compression_useful_column_all_cuts}
    \lean{TNLean.Lattice.Compression.DistributedOperatorContraction.exists_useful_column_of_sub_pure_le}
    \leanok
    \uses{def:compression_operator_operator,def:compression_operator_alphabets,
        def:compression_operator_endpoint_tensor,def:peps_dependent_network,
        def:peps_ket_column_dependent_fix_input}
    Let $\sigma$ be the distributed operator contraction with
    input and output alphabet $X_a$ at each party, all finite. Retain its
    tagged link alphabets $D_e$ and endpoint tensors $A_a$.
    Let $\omega\in\bigotimes_a\mathbb C^{X_a}$ be a unit vector and suppose
    $\|\sigma-|\omega\rangle\langle\omega|\|_1\le\epsilon<1$.
    There exist one product-basis input $z$ and one nonzero column
    $v=\sigma|z\rangle$. With the same identity bonds, its coefficients are
    \begin{align}
        A_{a,z}(\eta,x_a) & =A_a(\eta;(x_a,z_a)),                &
        v(x)              & =\sigma(x,z)=\mathcal N_{A_z}(1)(x).
        \label{eq:compression_useful_column_tensor}
    \end{align}
    For $\widehat v=v/\|v\|$, one has $\|\widehat v\|=1$ and some real
    $\theta$ satisfies $\|\widehat v-e^{i\theta}\omega\|\le2\epsilon$.
    The same $z,v$ satisfy, for every party subset $S$,
    \begin{align}
        \operatorname{rank}F_S(v)
         & \le\prod_{e\in\partial S}|D_e|, &
        \operatorname{rank}F_S(\widehat v)
         & =\operatorname{rank}F_S(v).
        \label{eq:compression_useful_column_ranks}
    \end{align}
    The cut product contains one alphabet for each tagged crossing
    edge, retaining parallel links and allowing empty alphabets.
    The column selection is that of
    \cite[Lemma~8.1]{OpenAI2026PolynomialPEPS}.
\end{lemma}
\begin{proof}\leanok
    \uses{thm:peps_ket_column_dependent,thm:compression_operator_column_cut_rank}
    Apply column selection to the endpoint tensors
    of $\sigma$. Fixing $z_a$ locally gives this same column on the
    same graph and alphabets. The identity-bond factorization therefore
    bounds its rank at every cut. Since $v\ne0$, multiplication by the
    nonzero scalar $\|v\|^{-1}$ preserves each matrix rank, giving
    (\ref{eq:compression_useful_column_ranks}). The vector statement
    gives the unit norm and phase-error bound.
\end{proof}

\begin{lemma}[Uniform compression and useful block columns]
    \label{thm:chain_uniform_compressed_block_column}
    \lean{TNLean.MPS.AreaLaw.ChainSweep.exists_uniform_sampled_block_column}
    \leanok
    \uses{def:peps_distributed_construction,def:peps_distributed_construction_density,
        def:chain_party_placement,def:compression_operator_operator}
    Let $i$ range over an arbitrary index set $I$, with sizes $s_i\ge2$
    and natural lengths and widths $L_i,w_i$. Let $\mathcal C_i$ be given
    distributed constructions on the parties of these chains,
    with no messages in their recorded monomial expansions;
    write $W_i$ for their original circuits and $\rho_i$ for their
    unrenormalized physical densities. The index map $i\mapsto L_i$
    need not be injective.
    Fix natural $a,b,\chi$ with $a\le b$ and
    $8(2b(a-1))\le\chi$. Fix natural $n,d,r,u,v$ and real
    $C_N,C_D,C_K,C_S$, independent of $i$.
    Subtraction in $a-1$ is truncated at zero.
    For each member, all parties $j$ and all original nonprivate gates $g$
    satisfy the following resource bounds:
    \begin{align}
        N(W_i)                   & \le C_Ns_i^n, & d_{i,j}  & \le C_Ds_i^d,\notag \\
        |\operatorname{part}(g)| & \le a,        & n_j(W_i) & \le b.\notag
    \end{align}
    For each $i$ there is a natural number $K_i$ bounding the number of
    recorded monomials at every original nonprivate gate, with
    $K_i\le C_Ks_i^u$. Every recorded monomial has at most $r$ pair
    effects, and at every original nonprivate gate the moduli of the
    recorded coefficients sum to at most $C_Ss_i^v$.
    Let $\omega_i$ be unit vectors in the physical output spaces,
    and let $\eta_i\in\mathbb R$ satisfy
    \begin{align}
        \|\rho_i-|\omega_i\rangle\langle\omega_i|\|_1 & \le\eta_i.
        \label{eq:chain_uniform_pure_input}
    \end{align}
    Assume also that any two participants of one original nonprivate
    gate of $W_i$ have level distance at most two. These hypotheses
    concern every member of the family. For circuits constructed in
    Lemma~\ref{thm:chain_conditional_preparation}, the geometric
    conclusion follows from its derived stage support.

    For every real $p>0$, there exist $C_p\ge0$ and an integer $\alpha>0$,
    chosen before all $i$, such that every member has a tagged
    distributed operator network $\sigma_i$ satisfying
    \begin{align}
        \|\sigma_i-\rho_i\|_1 & \le s_i^{-p},\notag                                \\
        \|\sigma_i-|\omega_i\rangle\langle\omega_i|\|_1
                              & \le\epsilon_i:=s_i^{-p}+\eta_i,
        \label{eq:chain_uniform_compression_error}                                 \\
        1-\eta_i              & \le\operatorname{Re}\operatorname{Tr}\rho_i\le1, &
        |D_e|                 & \le C_ps_i^\alpha.
        \label{eq:chain_uniform_compression_dimensions}
    \end{align}
    The alphabet bound in (\ref{eq:chain_uniform_compression_dimensions})
    holds for every tagged link $e$.
    Every party meets at most $2b(a-1)$ links, and the two endpoints
    of each link participate in one common original nonprivate gate
    of this same $W_i$. Consequently every level prefix has at most
    $\chi$ tagged crossing links.
    Set $B_i=\max\{1,\lceil C_ps_i^\alpha\rceil_+\}$, a natural number.
    Every star or sample alphabet has size at most $B_i$.

    Only for a member with $\epsilon_i<1$, choose one product-basis
    input $z_i$ and a nonzero column $v_i=\sigma_i|z_i\rangle$.
    This is the same input-fixed endpoint network, with
    $\|\widehat v_i\|=1$ and
    $\|\widehat v_i-e^{i\theta_i}\omega_i\|\le2\epsilon_i$ for some real
    $\theta_i$. For every natural prefix index $t$, put
    $S_t=\{j:\ell(j)<t\}$. The same column satisfies
    \begin{align}
        \operatorname{rank}F_{S_t}(v_i)
         & \le\prod_{e\in\partial S_t}|D_e|\le B_i^{\chi},
        \label{eq:chain_uniform_column_rank}               \\
        \operatorname{rank}F_{S_t}(\widehat v_i)
         & =\operatorname{rank}F_{S_t}(v_i)\le B_i^{\chi}.
        \label{eq:chain_uniform_normalized_rank}
    \end{align}
    The constants precede every family member, its network, the
    member's condition $\epsilon_i<1$, the useful column and all cuts.
    This is the conditional compression and column step of
    \cite[Theorem~5.2, Remark~5.3 and Lemma~8.1]{OpenAI2026PolynomialPEPS}.
\end{lemma}
\begin{proof}\leanok
    \uses{thm:peps_separate_arity_compression,
        thm:chain_tagged_link_congestion,thm:compression_useful_column_all_cuts}
    Apply Theorem~\ref{thm:peps_separate_arity_compression} to the
    whole family, retaining its sampled tensors and tagged alphabets.
    Its density estimates and trace bounds give
    (\ref{eq:chain_uniform_compression_error}) and
    (\ref{eq:chain_uniform_compression_dimensions}).
    The degree bound $2b(a-1)$ and the common-original-gate witnesses
    concern this same network. Lemma~\ref{thm:chain_tagged_link_congestion}
    gives at most $8(2b(a-1))\le\chi$ crossings for each prefix.
    For each link, round its uniform alphabet bound upwards to $B_i$.
    Since $B_i\ge1$, the crossing product is at most
    $B_i^{|\partial S_t|}\le B_i^{\chi}$.
    After fixing a member with $\epsilon_i<1$, apply
    Lemma~\ref{thm:compression_useful_column_all_cuts} once to its
    network. Its one choice of $z_i,v_i$ has the
    input-fixed tensors and works for all prefixes, giving
    (\ref{eq:chain_uniform_column_rank}) and
    (\ref{eq:chain_uniform_normalized_rank}).
\end{proof}

\begin{lemma}[Party prefixes and original-site coefficient matrices]
    \label{thm:chain_party_prefix_site_coordinates}
    \lean{TNLean.MPS.AreaLaw.ChainSweep.rank_partyPrefix_eq_cutMatrix,
        TNLean.MPS.AreaLaw.ChainSweep.rank_partyPrefix_last_eq_cutMatrix}
    \leanok
    \uses{def:chain_terminal_coordinates,def:chain_sweep_stage_assignments,
        def:chain_party_placement,def:chain_open_mps}
    Let $L,w,q$ be arbitrary natural numbers, with the same owner and
    block-coordinate definitions, including the zero-width convention.
    Let $\psi$ be any vector on the physical party factors
    $\bigotimes_a\mathbb C^{d_a}$, where $d_a=q^{|M_a|}$.
    Define the original-site vector exactly by
    \begin{align}
        \Phi(x) & =\psi(\chi(x)).
        \label{eq:chain_party_vector_pullback}
    \end{align}
    For every natural $a$, put $S_a=\{b:\ell(b)<a\}$. If $a\le m$, then
    \begin{align}
        \operatorname{rank}F_{S_a}(\psi)
         & =\operatorname{rank}M_{at}(\Phi).
        \label{eq:chain_party_prefix_rank}
    \end{align}
    The last prefix $S_{m+1}$ includes every party, including $R_m$, and
    \begin{align}
        \operatorname{rank}F_{S_{m+1}}(\psi)
         & =\operatorname{rank}M_L(\Phi).
        \label{eq:chain_party_last_rank}
    \end{align}
    Scalar physical factors remain in these products. The prefix $S_m$
    ends at site $mt$; the final remainder is included by $S_{m+1}$.
    These are the original block coordinates in
    \cite[Part~II, proof of Theorem~B]{TNLean2026OneDimensional}.
\end{lemma}
\begin{proof}\leanok
    \uses{def:chain_terminal_coordinates,def:chain_sweep_stage_assignments,
        def:chain_party_placement,def:chain_open_mps}
    For $a\le m$, the owner satisfies
    $\ell(\operatorname{owner}(x))<a$ exactly when $x<at$.
    On each side, restricting that owner map and applying the
    block encoding gives a bijection between site configurations and
    the corresponding party configurations, including empty owner fibres.
    Encoding after gluing the two site configurations is the same as
    gluing their two encoded party configurations. Thus $M_{at}(\Phi)$
    is $F_{S_a}(\psi)$ with its rows and columns separately reindexed
    by these bijections, giving (\ref{eq:chain_party_prefix_rank}).
    Every party level is less than $m+1$, so the same argument for all
    sites gives (\ref{eq:chain_party_last_rank}). No auxiliary coordinate
    is changed by these physical configuration bijections.
\end{proof}

\begin{lemma}[Rank bounds at every original-site cut]
    \label{thm:chain_site_ranks_of_party_prefixes}
    \lean{TNLean.MPS.AreaLaw.ChainSweep.rank_cutMatrix_le_of_partyPrefix_bounds}
    \leanok
    \uses{thm:chain_party_prefix_site_coordinates,def:chain_open_mps}
    In (\ref{eq:chain_party_vector_pullback}), suppose $q\ge1$, $w\ge1$,
    $L\ge60w$ and $R\in\mathbb N$ with $R\ge1$. Suppose the
    physical party tensor has the prefix bounds
    \begin{align}
        \operatorname{rank}F_{S_a}(\psi) & \le R
        \label{eq:chain_physical_prefix_hypothesis}
    \end{align}
    for every integer $1\le a\le m$. Then every original-site cut
    $0\le k\le L$ satisfies
    \begin{align}
        \operatorname{rank}M_k(\Phi) & \le q^{40w}R.
        \label{eq:chain_all_site_rank_bound}
    \end{align}
    The vector $\psi$ is arbitrary; no norm or nonzero premise is needed.
\end{lemma}
\begin{proof}\leanok
    \uses{thm:chain_party_prefix_site_coordinates,thm:chain_cut_rank_growth,
        thm:chain_terminal_block_sizes}
    The coordinate equalities give the same bound $R$ at cuts $at$,
    for $1\le a\le m$. The endpoint cuts at $0,L$ have rank at most
    one, hence at most $R$. A cut below $mt$ lies between two adjacent
    completed-block boundaries, separated by $t=20w$. A cut at or beyond
    $mt$ lies in the final interval $[mt,L]$, of length less than
    $4t=80w$. Apply Lemma~\ref{thm:chain_cut_rank_growth} using the
    rank bound $R$ at both ends of the appropriate interval. This gives
    (\ref{eq:chain_all_site_rank_bound}). If $m=0$, the prefix assumption
    is empty and this final interval is the entire chain.
\end{proof}

\section{Physical information estimates and the chain preparation}

Use the chain Hamiltonians and gapped ground vectors of
Definition~\ref{def:chain_model}, the corrected one-band assignments of
Definition~\ref{def:chain_one_band_assignments}, and the labelled stage
history of Definition~\ref{def:chain_sweep_frame_path}.
All intervals below are half-open lattice intervals; the real exchange
window remains open. The copy labels, their order and their retained
owners are those of
\cite{gap:tnlean_one_dimensional_band_sweep_corrections}.
Write $I_\Omega(T:E)$ for mutual information in the same original
chain vector $\Omega$ throughout. For a circuit $C$, write $N(C)$ for
its number of nonprivate occurrences, $n_p(C)$ for the number involving
$p$, and $\operatorname{part}(g)$ for a gate's participant set.
The prescribed placement is $\ell(R_i)=i$ and $\ell(P_{j+1})=j$.

\begin{lemma}[Information regions of an interior band]
    \label{thm:chain_band_physical_information}
    \lean{TNLean.MPS.AreaLaw.OneBand.primitive_separation,
        TNLean.MPS.AreaLaw.OneBand.primitive_information_of_logarithmic_separation}
    \leanok
    \uses{def:chain_one_band_assignments,def:chain_model}
    Let $L,a,w,q$ be natural numbers, $w\ge1$ and $8w\le a$.
    Let the original main owner be $C$ on
    $\lbrack a-8w,a)_L$, $A$ on $\lbrack a,a+w)_L$ and $B$ on
    $[a+w,a+2w)_L$, with $A\ne C$ and $B\ne C$.
    No owner is prescribed outside these intervals, and $A=B$ is allowed.
    For all eight pairs $(T_j,E_j)$ in the table of
    Lemma~\ref{thm:chain_one_band_conditional_changes},
    \begin{align}
        x\in T_j,\ y\in E_j & \quad\Longrightarrow\quad |x-y|\ge w.
        \label{eq:chain_band_separation}
    \end{align}
    In particular the exchange uses $E=Z\setminus Y$ and not the
    complement of $Y$; the death rows use their after-death sheets.

    Fix $q\ge1$, $R\in\mathbb N$, $J\in\mathbb R$ and $\Delta>0$.
    There are $D_{24}>0$ and $L_0\ge1$, depending only on these fixed
    parameters, such that the following holds for every $L\ge L_0$,
    every Hamiltonian of range $R$ and strength bound $J$ with gapped
    ground vector $\Omega$ and energy $E_0$, every choice of parties,
    and all the ownership data just stated.
    If $D_{24}\log L\le w$, then
    \begin{align}
        I_\Omega(T_j:E_j) & \le L^{-24},\qquad 0\le j<8.
        \label{eq:chain_band_information}
    \end{align}
    The choice of $D_{24},L_0$ precedes all chain, vector, width and
    ownership data. This is the logarithmic-separation argument for
    \cite[Part~II, the band dissolution lemma]{TNLean2026OneDimensional},
    with the stated lattice convention and exponent $24$.
\end{lemma}
\begin{proof}
    \leanok
    \uses{thm:chain_one_band_geometry,thm:chain_logarithmic_separation}
    In a sheet, $E(F,P)$ consists exactly of sites not
    owned by $P$. After the first main death the owner is $A$ on
    $[a-5w,a+w)_L$, and after the second it is $B$ on
    $[a-6w,a+2w)_L$. The left parts of these intervals inherit the
    auxiliary owner through the open exchange window; the right parts
    use only the stated $A$ and $B$ flanks. These two complete owner
    windows give the margins for the next three moves. The auxiliary
    births give their margins directly, and the first has empty
    environment. The exact exchange regions in
    Lemma~\ref{thm:chain_one_band_geometry} give its margin; the main
    endpoint gives the last margin. This proves
    (\ref{eq:chain_band_separation}).
    Each $T_j$ is one lattice interval, including the exchange
    target. Apply the logarithmic-separation estimate of
    \cite[Part~II, the lemma on information at logarithmic separation]{TNLean2026OneDimensional}
    with exponent $24$ and one interval. It chooses uniform constants
    before any of the remaining data and proves
    (\ref{eq:chain_band_information}) for the original $\Omega$.
\end{proof}

\begin{lemma}[Information for the two births and the first death]
    \label{thm:chain_stage_physical_information}
    \lean{TNLean.MPS.AreaLaw.ChainSweep.birth_death_separation,
        TNLean.MPS.AreaLaw.ChainSweep.birth_death_information_of_logarithmic_separation}
    \leanok
    \uses{def:chain_sweep_stage_assignments,def:chain_model}
    For arbitrary natural $L,q$, let $w\ge1$ and $0\le k<m$.
    Put $a=kt$ and $b=(k+1)t$, where $t=20w$.
    The target and environment pairs are
    \begin{align}
        (T_k^R,E(M_k,R_k))
         & =(\lbrack b+w,L)_L,\lbrack 0,a)_L),\notag               \\
        (T_k^P,E(M_k^{(1)},R_k))
         & =(\lbrack a+w,b)_L,\lbrack 0,a)_L\cup\lbrack b+w,L)_L).
        \label{eq:chain_birth_regions}
    \end{align}
    At $k=0$, the direct death has target $[0,w)_L$ and
    after-death environment $E(M_0^{(L)},P_1)=[t,L)_L$.
    Every such target is separated from its stated environment by at
    least $w$.

    For fixed $q\ge1$, $R\in\mathbb N$, $J\in\mathbb R$ and
    $\Delta>0$, there are $D_{24}>0$ and $L_0\ge1$ before all chain,
    Hamiltonian, ground-vector, width and stage data, such
    that $L\ge L_0$ and $w\ge D_{24}\log L$, with $w\ge1$, imply
    $I_\Omega(T:E)\le L^{-24}$ for both birth pairs and, at $k=0$,
    for the direct-death pair. Here $\Omega$ is any gapped ground
    vector of the given Hamiltonian. These are the three main-sheet
    estimates in \cite[Part~II, proof of
        the sweep proposition]{TNLean2026OneDimensional}.
\end{lemma}
\begin{proof}
    \leanok
    \uses{thm:chain_sweep_stage_owners,thm:chain_logarithmic_separation}
    The stage ownership gives (\ref{eq:chain_birth_regions})
    and the first after-death environment. Their interval endpoints
    give $|x-y|\ge w$ directly. Apply the same physical
    logarithmic-separation estimate with exponent $24$ and one
    interval to each target. No estimate is applied to a vector
    produced during the sweep.
\end{proof}

\begin{lemma}[Local circuits from physical logarithmic separation]
    \label{thm:chain_physical_local_circuits}
    \lean{TNLean.MPS.AreaLaw.ChainSweep.local_stage_circuits_of_logarithmic_separation}
    \leanok
    \uses{def:chain_model,def:chain_stage_participants,def:chain_sweep_frame_path,
        def:lattice_bounded_change,def:peps_circuit_source_bound}
    Fix $q\ge1$, $R\in\mathbb N$, $J\in\mathbb R$ and $\Delta>0$.
    There are $D_{24}>0$ and $L_0\ge1$ such that, for every
    $L\ge L_0$, every corresponding gapped ground vector $\Omega$
    and every natural $w\ge1$ with
    $D_{24}\log L\le w$, each stage $k<m$ admits eighteen
    circuits between its prescribed successive boundaries.
    For each $0\le j<18$, $c_{k,j}$ is a bounded change on
    $\mathcal S_{k,j}$ with bounds $r=K=1$, pair-source bound one,
    at most one nonprivate occurrence, and reference error at most
    $4L^{-12}$. At $k=0$ the positions $3,\ldots,9$ are the
    identity circuits, have zero error and have no nonprivate occurrence.
    Thus these circuits satisfy the entire local conclusion of
    Lemma~\ref{thm:chain_conditional_local_circuits}, including the
    original copy labels and spectator order. No information estimate
    or circuit is an additional physical premise. This gives the local
    circuits of \cite[Part~II, the sweep construction
        and the sweep proposition]{TNLean2026OneDimensional}.
\end{lemma}
\begin{proof}
    \leanok
    \uses{thm:chain_band_physical_information,thm:chain_stage_physical_information,
        thm:chain_stage_margins,thm:chain_sweep_stage_owners,
        thm:chain_conditional_local_circuits}
    Take the maximum of the two information constants and the maximum
    of their length cutoffs. Since $\log L\ge0$, this common width
    meets both original conditions. The existence of $k<m$ forces
    $4t\le L$ and hence the required margin $3t\le L$.
    The flanks permit the band estimate at
    $(a,A,B,C)=(kt,R_k,P_{k+1},P_k)$ only for $k>0$, and at
    $((k+1)t,R_k,R_{k+1},P_{k+1})$ at every stage.
    Rewrite the right band's main entrance by its endpoint
    identity. Together with the two births and the exceptional first
    death, these supply precisely the nineteen information
    hypotheses of Lemma~\ref{thm:chain_conditional_local_circuits}.
    The ground-state hypothesis supplies $\|\Omega\|=1$.
    The resulting circuits retain the seven identity positions; no
    predecessor $P_0$ or left copy is introduced at stage zero.
\end{proof}

\begin{lemma}[Uniform physical preparation by the sweep]
    \label{thm:chain_physical_preparation}
    \lean{TNLean.MPS.AreaLaw.ChainSweep.exists_preparation_of_gappedGroundState}
    \leanok
    \uses{def:chain_model,def:chain_terminal_reference,def:chain_terminal_coordinates,
        def:chain_stage_participants,def:chain_party_placement}
    Fix $q\ge1$, $R\in\mathbb N$, $J\in\mathbb R$ and $\Delta>0$.
    There are $D\ge1$, $K>0$, $s\in\mathbb N$ and $L_0\ge1$,
    before all further physical data, with the following property.
    For every $L\ge L_0$, Hamiltonian of these parameters, gapped
    ground vector $\Omega$ with energy $E_0$,
    put $w=\lceil D\log L\rceil_+$ and $N=18m$.
    The constructed initial party vectors $x_p$ have norm one.
    Their product input $X=\bigotimes_p x_p$ admits a circuit $C$
    to the specified terminal physical factors and all retained
    private auxiliary registers, with output $\Psi=CX$.
    For every party $p$ the same $C$ has the stated lifetime bound.
    Every original nonprivate gate $g$ has a stage $k<m$ and position
    $j<18$ such that
    \begin{align}
        N(C)                     & \le18m\le20m\le2L,         & n_p(C) & \le40, & \|\Psi\| & \le1,\notag \\
        |\operatorname{part}(g)| & \le3,                      &
        \operatorname{part}(g)   & \subseteq\mathcal S_{k,j}.
        \label{eq:chain_physical_resources}
    \end{align}
    Every original nonprivate gate has at most one recorded monomial,
    absolute coefficient sum at most one, and at most one normalized
    pair effect and one normalized pair source per monomial.
    In particular its arity is also at most $40$, and any two
    participants have level distance at most two.
    The geometric and dimension conclusions are
    \begin{align}
        1\le w,       & \qquad 60w\le L, & |\mathcal P| & \le2L,\notag \\
        d_p=q^{|M_p|} & \le q^{80w},     & d_p          & \le KL^s.
        \label{eq:chain_physical_dimensions}
    \end{align}
    Use the inverse terminal regrouping to form $Y$, the
    sheet readout $\mathcal R$, and the standard physical/private
    readout $V$ as in Lemma~\ref{thm:chain_conditional_preparation}.
    Put $\rho=\operatorname{Tr}_{\rm private}|V\rangle\langle V|$
    and $\varepsilon=4NL^{-12}$. Then
    \begin{align}
        \|\Xi\|                          & =1,                               &
        \|\mathcal RY-\Omega\otimes\Xi\| & \le\varepsilon\le26L^{-11},\notag   \\
        \|\rho-|U_\chi\Omega\rangle\langle U_\chi\Omega|\|_1
                                         & \le2\varepsilon\le52L^{-11}.
        \label{eq:chain_physical_errors}
    \end{align}
    The vector $\Xi$ is the product over the distinct original copies,
    including the empty product. It need not be a product over parties.
    Only the main factors have dimensions $d_p$; all auxiliary registers
    remain private at their owners. This is the preparation
    assertion of \cite[Part~II, the sweep proposition and
        the density estimate in the proof of Theorem~B]{TNLean2026OneDimensional},
    with \cite{gap:tnlean_one_dimensional_band_sweep_corrections}.
\end{lemma}
\begin{proof}
    \leanok
    \uses{thm:chain_physical_local_circuits,thm:chain_conditional_preparation,
        thm:chain_stage_arity_three,thm:chain_terminal_block_sizes}
    Enlarge the physical local constant to $D=\max\{D_{24},1\}$.
    Enlarge the length cutoff to at least $120$ and the uniform
    logarithmic threshold at which $60D\log L\le L/2$.
    Then $w<D\log L+1$ gives $60w\le L$, and $D\log L\ge1$
    gives $w\ge1$. One may take $K=q^{80}$ and
    $s=\lceil80D\log q\rceil_+$, since
    \begin{align}
        q^{80w} & \le q^{80}L^{80D\log q}\le KL^s.
        \label{eq:chain_physical_ceiling}
    \end{align}
    This retains the ceiling correction also at $q=1$.
    Choose the physical local circuits at every stage and apply
    Lemma~\ref{thm:chain_conditional_preparation}. This constructs
    the private initialization, chronological middle and final private
    regrouping with the same output and marginal. Its derived stage
    support, together with Lemma~\ref{thm:chain_stage_arity_three},
    gives arity three. The block geometry gives the party, count and
    dimension bounds. Since $18m\le L$,
    $\varepsilon\le4L^{-11}\le26L^{-11}$; the marginal bound
    is at most $2\varepsilon$. No normalization of $\Psi$ is made.
\end{proof}

\begin{definition}[The distributed construction with its original input]
    \label{def:chain_physical_construction}
    \lean{TNLean.MPS.AreaLaw.ChainSweep.preparationConstruction}
    \leanok
    \uses{def:peps_distributed_construction,thm:chain_initial_copies,
        def:chain_terminal_coordinates}
    Let $L>0$, $q\ge1$ and $w\in\mathbb N$.
    Given an original vector $\Omega$, the initial
    party spaces $H_p$ and their unit vectors $x_p$ from the grouped
    initial list, and a circuit $C$ from their product layout to the
    specified terminal layout, take the distributed construction
    with initial spaces $H_p$, initial vectors $x_p$, circuit $C$,
    physical dimensions $q^{|M_p|}$, the prescribed party order, and
    the full private auxiliary layout.
    The input is exactly $X=\bigotimes_p x_p$ and its physical density
    is the standard readout of $CX$ followed by the partial trace.
    The party $R_0$ exists even if $m=0$; parties with scalar main
    outputs and the final remainder $R_m$ are retained.
    The definition imposes no resource or approximation bound on $C$.
    The input and output factors follow
    \cite[Part~II, the sweep proposition]{TNLean2026OneDimensional},
    with \cite{gap:tnlean_one_dimensional_band_sweep_corrections}.
\end{definition}

\begin{lemma}[A preparation family over all physical instances]
    \label{thm:chain_uniform_physical_family}
    \lean{TNLean.MPS.AreaLaw.ChainSweep.exists_uniform_preparation_family}
    \leanok
    \uses{def:chain_physical_construction,def:chain_model,
        def:chain_terminal_coordinates}
    For fixed $q\ge1$, $R\in\mathbb N$, $J\in\mathbb R$ and
    $\Delta>0$, there are $D\ge1$, $K>0$, $s\in\mathbb N$ and
    $L_0\ge2$ before every physical instance.
    Let $\iota$ consist of all tuples $(L,H,\Omega,E_0,\mathrm{pos})$
    with $L\ge L_0$, $H$ a Hamiltonian of the fixed parameters and
    $\Omega$ a gapped ground vector at energy $E_0$.
    For every $i\in\iota$, set $w_i=\lceil D\log L_i\rceil_+$.
    One can choose the unit initial vectors and the circuit
    $C_i$ together, obtaining the construction of
    Definition~\ref{def:chain_physical_construction} for all $i$.
    With $\omega_i=U_{\chi_i}\Omega_i$, these same constructions satisfy
    \begin{align}
        L_i                      & \ge2,           & N(C_i)    & \le2L_i,            & d_{i,p} & \le KL_i^s,\notag \\
        |\operatorname{part}(g)| & \le3,           & n_p(C_i)  & \le40,              &
        \|\omega_i\|             & =1,\notag                                                                       \\
        \|\rho_i-|\omega_i\rangle\langle\omega_i|\|_1
                                 & \le52L_i^{-11}, & 1\le w_i, & \qquad60w_i\le L_i.
        \label{eq:chain_uniform_physical_family}
    \end{align}
    Every gate has monomial bound one, every recorded monomial has at
    most one pair effect, and the coefficient moduli at each gate sum to
    at most one. Its original participants
    have pairwise level distance at most two.
    The target density is the coordinate reindexing of
    $|\Omega_i\rangle\langle\Omega_i|$ by $\chi_i$.
    Thus all premises of Lemma~\ref{thm:chain_uniform_compressed_block_column}
    hold with arity $3$, lifetime $40$, crossing parameter $1600$,
    size $L_i$, gate-count parameters $(n,C_N)=(1,2)$,
    physical-dimension parameters $(d,C_D)=(s,K)$,
    $r=1$, $u=v=0$, $C_K=C_S=1$ and $\eta_i=52L_i^{-11}$,
    using the uniform compression theorem of
    \cite[Theorem~5.2]{OpenAI2026PolynomialPEPS}.
    No one Hamiltonian or ground vector is selected to represent a length.
\end{lemma}
\begin{proof}\leanok
    \uses{thm:chain_physical_preparation,def:chain_physical_construction,
        def:peps_distributed_construction_density}
    Enlarge the physical cutoff to $\max\{L_0,2\}$ and choose the
    preparation witnesses over the entire admissible tuple set.
    Retain the derived arity-three bound of these same circuits.
    Record them without changing the circuit, input, private layout
    or standard marginal. The pair-effect and coefficient-sum bounds
    are the expansion bound $(1,1)$ of the preparation. The main-coordinate map is an isometry,
    and its pure density equals the reindexing entry by entry.
    Hence the physical preparation estimates give
    (\ref{eq:chain_uniform_physical_family}) and every stated
    compression premise for this same family.
\end{proof}

\section{Polynomial bounds and the physical approximation}

\begin{lemma}[The two ceiling corrections in the bond bound]
    \label{thm:chain_protocol_polynomial_bond}
    \lean{TNLean.MPS.AreaLaw.ChainSweep.polynomial_bond_bound_of_ceil_width}
    \leanok
    Let $q\ge1$, $D\ge1$, $C_p\ge0$, a natural $\chi$, and $\alpha\in\mathbb N$
    with $\alpha>0$. For $L\ge2$ and $w=\lceil D\log L\rceil_+$,
    put $A_L=\max\{1,\lceil C_pL^\alpha\rceil_+\}$ and
    $n=\max\{1,\lceil40D\log q\rceil_++\chi\alpha\}$.
    Then
    \begin{align}
        q^{40w}A_L^{\chi}
         & \le q^{40}(C_p+1)^{\chi}L^n.
        \label{eq:chain_protocol_polynomial_bond}
    \end{align}
    Both the exponent and the prefactor are independent of $L$.
    The inequality includes $q=1$, $C_p=0$ and $\chi=0$.
\end{lemma}
\begin{proof}\leanok
    The ceiling estimates give
    $A_L\le(C_p+1)L^\alpha$ and
    $w\le D\log L+1$.
    Nonnegativity of $\log q$ and $\log L$ yields
    $q^{40w}\le q^{40}L^{\lceil40D\log q\rceil_+}$.
    Raise the alphabet estimate to the natural power $\chi$,
    multiply the bounds and increase the exponent to $n$.
    If $\chi=0$, the alphabet factor is $1$ and $n\ge1$ still holds.
    No division by $\log q$ occurs. This is the polynomial estimate
    in \cite[Part~II, end of the proof of Theorem~B]{TNLean2026OneDimensional},
    retaining both additive ceiling corrections.
\end{proof}

\begin{lemma}[An MPS of the same physical vector in site coordinates]
    \label{thm:chain_protocol_mps_of_prefix_ranks}
    \lean{TNLean.MPS.AreaLaw.ChainSweep.isOpenMPS_of_polynomial_partyPrefix_bounds}
    \leanok
    \uses{def:chain_open_mps,def:chain_terminal_coordinates,
        def:chain_party_placement}
    Retain $q,D,C_p,\alpha,\chi,L,w,A_L,n$ from
    Lemma~\ref{thm:chain_protocol_polynomial_bond} and assume
    $60w\le L$. Let $\psi$ be a unit vector in the physical
    party space. Suppose its coefficient flattening satisfies
    $\operatorname{rank}F_{S_r}(\psi)\le A_L^{\chi}$ for every
    $1\le r\le m$, where $S_r=\{p:\ell(p)<r\}$.
    Write $\chi_{\rm cfg}$ for the fixed main-configuration equivalence.
    Define the inverse main-coordinate vector
    $\Phi(x)=\psi(\chi_{\rm cfg}(x))$ and $B=q^{40w}A_L^{\chi}$.
    Then $\|\Phi\|=1$, every original-site cut $0\le k\le L$
    satisfies $\operatorname{rank}M_k(\Phi)\le B$, and $\Phi$
    is an open-boundary MPS of bond $B$. The polynomial bound
    (\ref{eq:chain_protocol_polynomial_bond}) holds and $n\ge1$.
    The prefix-rank bounds are hypotheses here; for the protocol vector
    they are proved in Theorem~\ref{thm:chain_protocol_mps_approximation}.
    It is the block-to-site passage in
    \cite[Part~II, end of the proof of Theorem~B]{TNLean2026OneDimensional}.
\end{lemma}
\begin{proof}\leanok
    \uses{thm:chain_site_ranks_of_party_prefixes,thm:chain_mps_of_cut_rank,
        thm:chain_protocol_polynomial_bond}
    Positivity of $D\log L$ gives $w\ge1$, and $A_L^{\chi}\ge1$.
    The coordinate equivalence preserves the norm.
    Apply Lemma~\ref{thm:chain_site_ranks_of_party_prefixes} with
    $R=A_L^{\chi}$. It retains scalar physical parties and the
    nonempty final remainder, and bounds every original-site cut by
    $q^{40w}R$, including when $m=0$.
    Lemma~\ref{thm:chain_mps_of_cut_rank} constructs an MPS
    representation of this very $\Phi$. Apply
    Lemma~\ref{thm:chain_protocol_polynomial_bond} to its bond.
\end{proof}

\begin{theorem}[Polynomial MPS approximation by the physical protocol]
    \label{thm:chain_protocol_mps_approximation}
    \lean{TNLean.MPS.AreaLaw.ChainSweep.exists_protocol_isOpenMPS_norm_sub_phase_le}
    \leanok
    \uses{def:chain_model,def:chain_open_mps}
    Fix $q\in\mathbb N$ with $q\ge1$, $R\in\mathbb N$,
    $J\in\mathbb R$ and $\Delta>0$. There are $C>0$ and
    $n\in\mathbb N$ with $n>0$, depending only on these fixed
    parameters, such that the following holds for every $L\ge2$,
    every Hamiltonian of range $R$ and strength bound $J$, and every
    gapped ground vector $\Omega$ at energy $E_0$.
    There are $B\in\mathbb N$ and a vector $\Phi$ on the original
    chain with
    \begin{align}
        \|\Phi\| & =1, & \Phi & \ne0, & B & \le CL^n,
        \label{eq:chain_protocol_headline_bond}
    \end{align}
    such that $\Phi$ is an open-boundary MPS of bond $B$, and some
    $\theta\in\mathbb R$ satisfies
    \begin{align}
        \|\Phi-e^{i\theta}\Omega\| & \le L^{-1}.
        \label{eq:chain_protocol_headline_error}
    \end{align}
    This is \cite[Theorem~B]{TNLean2026OneDimensional}, with a unit
    approximant and a positive natural polynomial exponent. Neither
    information estimates, circuit data nor rank bounds are hypotheses
    of this theorem.
\end{theorem}
\begin{proof}\leanok
    \uses{thm:chain_uniform_physical_family,
        thm:chain_uniform_compressed_block_column,
        thm:chain_protocol_mps_of_prefix_ranks,thm:chain_mps_of_cut_rank}
    Take the entire admissible physical family from
    Lemma~\ref{thm:chain_uniform_physical_family}, with its derived
    arity-three bound. Apply Lemma~\ref{thm:chain_uniform_compressed_block_column}
    to this whole family with $a=3$, $b=40$, $\chi=1600$ and
    accuracy exponent $5$. The numerical conditions follow from
    $3\le40$ and $8(2\cdot40(3-1))=1280\le1600$.
    Its uniform $C_p\ge0$ and natural $\alpha>0$ precede every
    member. Its own sampled network has degree at most $160$,
    common-original-gate witnesses, at most $1280$ crossings, and
    the resulting upper bound $1600$ used for its column ranks.
    Fix
    \begin{align}
        n             & =\max\{1,\lceil40D\log q\rceil_++1600\alpha\},\notag   \\
        C_{\rm large} & =q^{40}(C_p+1)^{1600},                               &
        C             & =\max\{C_{\rm large},q^{L_0}\}.
        \label{eq:chain_protocol_uniform_constants}
    \end{align}
    These constants precede all $L,H,\Omega,E_0$.

    For $L\ge L_0$, select the member of the family given by the chain,
    Hamiltonian, ground vector and energy. Put
    $\epsilon=L^{-5}+52L^{-11}$. For every $L\ge2$,
    $52L^{-11}\le L^{-5}$, $\epsilon<1$ and
    $2\epsilon\le4L^{-5}\le L^{-1}$.
    The operator network therefore has one nonzero
    product-basis column $v$, chosen before all cuts, whose
    normalization $u=v/\|v\|$ has unit norm and phase error at most
    $2\epsilon$ from $U_{\chi_{\rm cfg}}\Omega$, where
    $\chi_{\rm cfg}$ is the main-configuration equivalence.
    Fixing the physical input changes only the local tensors; the
    same graph and its tagged alphabets remain. Every alphabet has
    dimension at most $A_L=\max\{1,\lceil C_pL^\alpha\rceil_+\}$,
    so this same $u$ has all prefix ranks at most $A_L^{1600}$.
    Normalization does not change those ranks.
    Apply Lemma~\ref{thm:chain_protocol_mps_of_prefix_ranks} to
    $u$. It gives the original-site vector
    $\Phi(x)=u(\chi_{\rm cfg}(x))$ with bond $B=q^{40w}A_L^{1600}$ and
    $B\le C_{\rm large}L^n\le CL^n$.
    The coordinate isometry gives, for the same phase,
    \begin{align}
        \|\Phi-e^{i\theta}\Omega\|
         & =\|u-e^{i\theta}U_{\chi_{\rm cfg}}\Omega\|
        \le2\epsilon\le L^{-1}.\notag
    \end{align}

    If $2\le L<L_0$, instead take $\Phi=\Omega$ and $B=q^{L_0}$.
    Each cut rank is at most its row-configuration cardinality,
    hence at most $q^{L_0}$. The internal-cut rank theorem gives
    an exact MPS representation. Its phase error is zero, and
    $B\le C\le CL^n$. In both cases norm one implies $\Phi\ne0$.
    Thus the single uniform choice
    (\ref{eq:chain_protocol_uniform_constants}) proves
    (\ref{eq:chain_protocol_headline_bond}) and
    (\ref{eq:chain_protocol_headline_error}).
\end{proof}

\subsection{Sharper crossings for the sampled physical network}

\begin{lemma}[Crossing links for a three-party chain preparation]
    \label{thm:chain_sharp_sampled_crossings}
    \lean{TNLean.MPS.AreaLaw.ChainSweep.exists_uniform_sampled_network_with_prefixCrossingLinks_le_1600}
    \leanok
    \uses{def:chain_party_placement,def:peps_distributed_construction}
    Consider a family on the chain parties with the polynomial count,
    physical-dimension, monomial and coefficient-sum hypotheses of
    Theorem~\ref{thm:peps_separate_arity_compression}. Suppose its
    original circuits have arity at most three and lifetime at most
    forty. Let $\omega_i$ be unit physical vectors with
    $\|\rho_i-|\omega_i\rangle\langle\omega_i|\|_1\le\eta_i$.
    Suppose any two participants of one original nonprivate gate
    have placement distance at most two.
    The real size $L_i\ge2$ may be independent of the natural chain
    length and width used to specify the party set.

    For each $p>0$ there are $C_p\ge0$ and natural $\alpha>0$,
    before all members, and for every member one sampled
    network with both density estimates, retained trace bounds,
    polynomial alphabets and common-original-gate witnesses in
    (\ref{eq:peps_separate_arity_network})--(\ref{eq:peps_separate_arity_pure}).
    On this same tagged network every party has degree at most
    $160$, and every natural prefix index $r$ satisfies
    \begin{align}
        |\partial_r| & \le1280\le1600.
        \label{eq:chain_sharp_sampled_crossings}
    \end{align}
    Here $\partial_r$ consists of the links crossing
    $\{j:\ell(j)<r\}$. These are auxiliary preparation hypotheses
    for the crossing step in
    \cite[Part~II, proof of Theorem~B]{TNLean2026OneDimensional}.
\end{lemma}
\begin{proof}\leanok
    \uses{thm:peps_separate_arity_compression,thm:chain_tagged_link_congestion}
    Apply Theorem~\ref{thm:peps_separate_arity_compression} with
    $a=3$ and $b=40$. It retains the network and gives
    $2\cdot40(3-1)=160$ incident links per party.
    Its common-original-gate witnesses and the original gate's
    distance bound put the endpoints of each crossing link in
    $\{j:\max\{r-2,0\}\le\ell(j)<r+2\}$.
    There are at most four levels and two party types at each level.
    Lemma~\ref{thm:chain_tagged_link_congestion} therefore gives
    $|\partial_r|\le8\cdot160=1280$.
    This counts the original tags, including parallel links, and
    retains scalar-output parties and the final remainder $R_m$.
    The empty endpoint prefixes and zero-stage cases are included.
\end{proof}

\begin{lemma}[The sharper network bound from physical chain data]
    \label{thm:chain_physical_sharp_sampled_crossings}
    \lean{TNLean.MPS.AreaLaw.ChainSweep.exists_physical_sampled_network_with_prefixCrossingLinks_le_1600}
    \leanok
    \uses{def:chain_model,def:chain_physical_construction,def:chain_party_placement}
    Fix $q\in\mathbb N$ with $q\ge1$, $R\in\mathbb N$,
    $J\in\mathbb R$ and $\Delta>0$.
    There are $D\ge1$, $K>0$, $s\in\mathbb N$ and $L_0\ge2$
    before all physical instances. Let $\iota$ be the entire set
    of admissible tuples $(L,H,\Omega,E_0,\mathrm{pos})$ with
    $L\ge L_0$, Hamiltonian parameters $q,R,J$, and gapped
    ground vector $\Omega$ of gap $\Delta$ at energy $E_0$.
    Set $w_i=\lceil D\log L_i\rceil_+$.
    There are unit initial party vectors and circuits $C_i$ for
    all $i\in\iota$ giving the constructions of
    Definition~\ref{def:chain_physical_construction} and retaining
    every resource, target-density and width conclusion of
    Lemma~\ref{thm:chain_uniform_physical_family}.
    Their arity is at most three.
    In particular, for the same original circuits,
    \begin{align}
        N(C_i)                                        & \le2L_i,        & d_{i,j} & \le KL_i^s, & n_j(C_i) & \le40,\notag \\
        |\operatorname{part}(g)|                      & \le3,           &
        \|\rho_i-|\omega_i\rangle\langle\omega_i|\|_1 & \le52L_i^{-11},
        \label{eq:chain_sharp_physical_inputs}
    \end{align}
    where $\omega_i=U_{\chi_i}\Omega_i$ is the unit
    main-coordinate vector. Monomial and coefficient-sum bounds are one,
    original-gate participant distances are at most two, and
    $w_i\ge1$, $60w_i\le L_i$.

    There are $C_p\ge0$ and natural $\alpha>0$, before every
    member of this entire family, such that each member has one
    sampled network $\sigma_i$ satisfying
    \begin{align}
        \|\sigma_i-\rho_i\|_1 & \le L_i^{-5},\notag                                                        \\
        \|\sigma_i-|\omega_i\rangle\langle\omega_i|\|_1
                              & \le L_i^{-5}+52L_i^{-11},\notag                                            \\
        1-52L_i^{-11}         & \le\operatorname{Re}\operatorname{Tr}\rho_i\le1,\notag                     \\
        |D_e|                 & \le C_pL_i^\alpha,                                     & |I_j| & \le160, &
        |\partial_r|          & \le1280\le1600.
        \label{eq:chain_sharp_physical_outputs}
    \end{align}
    The incidence and prefix bounds hold at every party and every
    natural prefix index of this same network. Every link's
    endpoints participate in a common original gate of $C_i$.
    The input, circuit, physical coordinate map and unnormalized
    density are those of the constructed preparation.
    Only the main registers are physical: every auxiliary copy
    remains private at its specified retained copy-site owner, as in
    \cite{gap:tnlean_one_dimensional_band_sweep_corrections}.
    No information, circuit, arity or crossing bound is an additional
    physical premise. This supplies the sampled-network crossing
    estimate in \cite[Part~II, proof of Theorem~B]{TNLean2026OneDimensional}.
\end{lemma}
\begin{proof}\leanok
    \uses{thm:chain_uniform_physical_family,thm:chain_sharp_sampled_crossings}
    In the physical preparation family retain the derived arity-three
    bound of the same chosen circuits. The arity follows from their
    stage support and the participant table; it is independent
    of the lifetime bound. All members are present before sampling
    constants are selected.
    Apply Lemma~\ref{thm:chain_sharp_sampled_crossings} at $p=5$,
    with count parameters $(n,C_N)=(1,2)$, physical-dimension
    parameters $(d,C_D)=(s,K)$, $r=1$, $u=v=0$,
    $C_K=C_S=1$ and $\eta_i=52L_i^{-11}$.
    Retain its witnesses and all original preparation conclusions.
    This gives (\ref{eq:chain_sharp_physical_outputs}) on the same
    tensor entries and tagged alphabets.
\end{proof}
